% !TeX root = ../../Book2.tex
% 纲要标题：## 第12章*　障碍理论与 Postnikov 分解
% 纲要位置：几何学计划纲要.md，第 1468 行。
% 正文状态：初稿；尚未正式审核。

% 纲要内容开始
% 
% > 本章目标与证明范围：完整建立简单空间或明确局部系数条件下的首障碍与逐层延拓，随后构造 Postnikov 阶段。障碍群、具体障碍类和消失后的选择自由分别处理。
% 
% 纲要内容结束

\OptionalChapter{障碍理论与 Postnikov 分解}{g2:ch:12}
\BookChapterNumber{g2:ch:12}

\begin{chapterabstract}
	沿胞腔延拓一张映射或一张截面时，边界留下的球面映射决定下一步能否完成。本章把这些局部条件组成带同伦群系数的障碍余圈，证明余圈条件、改变选择的余边界公式、首障碍的选择无关性及延拓之间的差异分类，并由此证明上同调的 Eilenberg–MacLane 可表性。随后构造单连通空间的 Postnikov 塔，以映射柱和相对 Hurewicz 定理将每一阶段恢复为一张上同调类的同伦纤维。最后将理论用于球丛、向量场和结构群约化，比较首截面障碍与 Euler 类，并完成低维底空间上的存在性与同伦分类计算。
\end{chapterabstract}

\begin{learninggoals}
	\begin{itemize}
		\item 将相对映射延拓和截面延拓写成逐胞腔问题，识别各步实际使用的同伦群与局部系数。
		\item 构造障碍余圈，证明改变上一层选择产生余边界，区分固定当前映射与允许修改后的延拓条件。
		\item 说明首障碍何时与选择无关，并用差异类描述延拓存在后的自由度。
		\item 构造 Postnikov 阶段及其 Eilenberg–MacLane 纤维，使用 \(k\)-不变量恢复阶段并判断提升。
		\item 计算球丛和低维向量丛的截面障碍，说明独立截面如何给出结构群约化。
	\end{itemize}
\end{learninggoals}

在圆周上给定 \(f(z)=z^k\)，什么时候能把它填成一张 \(D^2\to S^1\) 的映射？答案是 \(k=0\)：圆盘中的同心圆可以把边界映射收缩到中心，所以一个非零绕数无法被填掉。若边界值来自某个丛中的单位向量，选取圆盘上的平凡化后，同样会出现一个绕数。于是“填一张映射”和“连续选一个向量”可以由同一种延拓判据处理。

\ProseParagraphBreak

把圆盘换成有许多胞腔的空间，问题多了一层：每张胞腔的边界可能单独绕出非零类，但修改相邻低维胞腔上的选择，也会改变这些类。障碍上链记录当前选择留下的全部边界；障碍上同调类则把这种修改带来的变化除去。这一区别正是“当前填不进去”与“调整上一层以后仍然填不进去”的区别，贯穿本章的存在性证明。

\ProseParagraphBreak

延拓还可以按目标的同伦次数来组织。一个空间的二次、三次同伦群分别是什么，并不能说明二次信息与三次信息怎样接在一起。Postnikov 阶段逐层保留这些群，\(k\)-不变量记录相邻阶段的粘合；本章将构造一族所有同伦群都相同、上同调却不同的空间，让这项额外信息直接出现在计算中。

\ProseParagraphBreak

本章使用第1章的相对 CW 同伦延拓性质、第5章的局部系数、第8章的纤维运输与同伦长正合列、第9章的 Hurewicz 定理和消去同伦群的构造。Postnikov 阶段与 Euler 类的比较还使用第11章的传递。读到相应位置时会重新指出这些结果所需的条件。首障碍与逐层延拓在简单纤维的范围内建立，非平凡的基底运输保留为局部系数；Postnikov 的重构证明在单连通空间中展开。几何应用涉及第一卷向量丛的标架、内积和正交补，所用构造在正文中作简要回顾。

% !TeX root = ../../Book2.tex
% 纲要标题：### 12.1 映射和截面的延拓问题
% 纲要位置：几何学计划纲要.md，第 1472 行。
% 正文状态：初稿；尚未正式审核。

% 纲要内容开始
% 
% * 空间对、骨架与相对条件
% * 单个胞腔上的延拓判据
% * 基本群作用产生的系数系统
% 
% 纲要内容结束

\section{映射和截面的延拓问题}
\label{g2:sec:1201}

给圆周上的每一点指定一个单位向量以后，能否把这个选择连续地填到整个圆盘？如果向量已经绕了一整圈，填充时便会遇到障碍。更高维的情形保留了同一结构：一个新胞腔是圆盘，旧数据已经规定它的边界；是否能够填进去，由边界在目标或纤维中的同伦类决定。先把这些逐胞腔的问题写准确，下一节才能把它们组织为上同调类。

\subsection{空间对、骨架与相对条件}
\label{g2:subsec:1201:01}

设 \((X,A)\) 为 CW 对，给定映射 \(f_A:A\to Y\)。\term[yingshe-yantuo-wenti]{映射延拓问题}[extension problem]要求寻找 \(f:X\to Y\)，满足 \(f|_A=f_A\)。这里边界数据是实际给定的映射；若只要求其同伦类相同，就需要另作说明。记
\begin{equation}\label{g2:eq:relative-obstruction-skeleta}
	X_{-1}=A,\qquad X_r=A\cup X^r\quad(r\ge0).
\end{equation}
这些是相对骨架：即使 \(A\) 含高维胞腔，它也从一开始就全部保留。下标 \(r\) 只控制 \(X\setminus A\) 中已经加入的胞腔；它与上一章表示同伦截段的 \(X[n]\) 含义不同。

\ProseParagraphBreak

更一般地，给定纤维化 \(p:E\to X\) 及 \(A\) 上的截面 \(s_A\)，要求寻找满足
\[
	p\circ s=\operatorname{id}_X,\qquad s|_A=s_A
\]
的截面。映射延拓是它的特例：取 \(E=X\times Y\)，截面必形如 \(s(x)=(x,f(x))\)。而对一般提升问题，若已有 \(g:X\to B\)、纤维化 \(p_0:E_0\to B\) 和 \(g|_A\) 的一个指定提升，则先取拉回
\[
	g^*E_0=\{(x,e):g(x)=p_0(e)\}\longrightarrow X.
\]
截面 \(x\mapsto(x,\widetilde g(x))\) 与满足原边界条件的提升 \(\widetilde g\) 一一对应。于是映射、提升和截面可以用同一套逐骨架方法处理。

\ProseParagraphBreak

从 \(X_r\) 到 \(X_{r+1}\) 时，每个新胞腔的边界已经有值，同维胞腔内部则互不相交。只要对每个特征圆盘找到保持边界的延拓，就能沿原来的附着映射粘合。CW 弱拓扑保证最后所得映射连续：连续性只需在 \(A\) 及每个特征圆盘上核对。对无限 CW 复形，这个事实允许逐维构造相容映射；在各个有限阶段分别找到互不相容的延拓，则还没有得到一张定义在 \(X\) 上的映射。

\begin{example}[圆盘边界上的绕数]\label{g2:exm:boundary-winding-extension}
	令 \(X=D^2\)、\(A=S^1\)，取 \(f_A(z)=z^k:S^1\to S^1\)，其中 \(k\in\mathbb Z\)。若它延拓到 \(D^2\)，把圆盘沿半径收缩到中心，就得到 \(f_A\) 的自由零伦，故绕数必须为零。反过来，\(k=0\) 时常值映射本身便给出延拓。因此延拓恰在 \(k=0\) 时存在。这同时是平凡实二维丛上的单位截面延拓问题：单位向量的边界绕数阻止它无零点地填入圆盘。
\end{example}

还须区分两种要求：把一个已经选好的 \(s_r:X_r\to E\) 原样延到下一层，和只固定 \(X_{r-1}\) 上的数据、允许重新选择 \(r\) 胞腔内的截面。第一种要求逐胞腔的障碍值全部为零；第二种允许用新选择改变这些值，因而将在下一节产生“障碍类为零”的判据。

\ProseParagraphBreak

比较两张延拓的同伦也属于同一问题。若 \(s_0,s_1\) 在 \(A\) 上相同，寻找固定 \(A\) 的截面同伦，就是在
\[
	\bigl(X\times I,\ (X\times\partial I)\cup(A\times I)\bigr)
\]
上延拓截面，其中两端取 \(s_0,s_1\)，侧面保持给定值。每个 \(r\) 胞腔乘区间成为一个 \((r+1)\) 维延拓问题，这将解释差异类的次数为何比存在性障碍低一度。

\subsection{单个胞腔上的延拓判据}
\label{g2:subsec:1201:02}

先看普通映射。设 \(q\ge1\)，\(\chi:D^{q+1}\to X_{q+1}\) 是一个新胞腔的特征映射，\(f_q:X_q\to Y\) 已经给定。复合
\[
	S^q\xrightarrow{\chi|_{S^q}}X_q\xrightarrow{f_q}Y
\]
就是待填圆盘上的边界数据。

\begin{proposition}[单个胞腔的映射延拓]\label{g2:prop:cell-extension-criterion}
	设 \(q\ge1\)，\(Y\) 为道路连通空间。映射 \(u:S^q\to Y\) 能延拓到 \(D^{q+1}\)，当且仅当它自由零伦；选择球面基点及目标基点识别后，这等价于对应的 \(\pi_q(Y)\) 类为零。因此，\(f_q:X_q\to Y\) 能保持已有值延到 \(X_{q+1}\)，当且仅当每个相对 \((q+1)\) 胞腔给出的这个类均为零。
\end{proposition}
\begin{proof}
	若 \(U:D^{q+1}\to Y\) 延拓 \(u\)，则 \(H(z,t)=U((1-t)z)\) 把边界映射收缩到 \(U(0)\)。反过来，给定从 \(u\) 到常值的同伦，将 \(D^{q+1}\) 写为 \(S^q\) 的锥，这个同伦就在锥顶处下降为所需延拓。

	由\cref{g2:cor:free-homotopy-orbits}，自由同伦对应基本群对带基点同伦群的作用轨道。任何作用自同构都保持零元，且只有零元能被送到零元，因此自由零伦等价于任一基点识别下的类为零。最后逐胞腔作上述填充；它们在公共边界均等于 \(f_q\)，所以可粘合为下一相对骨架上的连续映射。
\end{proof}

零胞腔要求目标非空；一胞腔要求它的两个端点像可以用道路相连。这两步是非空性和道路分支的问题。二胞腔开始才出现 \(\pi_1\) 中的边界环路；若这个群非交换，仍可判定单个环路是否为单位元，但尚不能把全部障碍当作普通阿贝尔群系数的上链。

\ProseParagraphBreak

截面的圆盘判据需要把边界数据送到同一纤维。设 \(p:E\to X\) 为 Hurewicz 纤维化，拉回到特征圆盘得到 \(\pi:V=\chi^*E\to D^{q+1}\)。第11章的\cref{g2:lem:fibration-disk-model}给出保持基底的纤维同伦等价
\[
	V\ \underset{P}{\stackrel{Q}{\rightleftarrows}}\ D^{q+1}\times F,
\]
其中 \(F=\pi^{-1}(0)\)。给定的边界截面 \(s_\partial:S^q\to V\) 经 \(Q\) 写成 \(Q s_\partial(z)=(z,u(z))\)。由此得到边界的纤维映射 \(u:S^q\to F\)。

\begin{proposition}[单个胞腔的截面延拓]\label{g2:prop:section-cell-extension}
	设 \(q\ge1\)，\(\pi:V\to D^{q+1}\) 为具有非空道路连通纤维的 Hurewicz 纤维化，\(s_\partial\) 为其边界上的截面。在上述纤维同伦乘积模型中，\(s_\partial\) 可延拓为整个圆盘上的截面，当且仅当 \(u:S^q\to F\) 的类为零。这个零值判据与乘积模型及基点识别的选择无关。
\end{proposition}
\begin{proof}
	若截面 \(s\) 已存在，\(Qs\) 的纤维坐标便延拓 \(u\)，所以其类为零。

	反过来，假定 \(u\) 延拓为 \(U:D^{q+1}\to F\)。复合 \(t(x)=P(x,U(x))\) 是圆盘上的截面；其边界未必等于 \(s_\partial\)，但 \(PQ\) 与恒等映射之间的纤维同伦给出从 \(t|_{S^q}\) 到 \(s_\partial\) 的截面同伦。对 CW 对 \((D^{q+1},S^q)\) 应用\cref{g2:lem:relative-homotopy-lift}，基底同伦取 \(H(x,r)=x\)，初始提升取 \(t\)，边界提升取刚才的同伦。所得提升的终时刻是一个严格延拓 \(s_\partial\) 的截面。这一步把纤维同伦模型给出的边界相等到同伦，修正成实际边界相等。

	判据两端已经等价于与模型无关的截面存在性，因此改变模型不改变它是否为零；基点改变也保持零元。
\end{proof}

纤维丛在圆盘上可用实际平凡化完成同样的论证；保留纤维同伦模型则同时适用于路径纤维化。对一般 Serre 纤维化，相同的圆盘论证可以只提升有限 CW 参数：沿径向收缩提升已给边界，得到中心纤维中的球面映射；若它零伦，就把这张圆盘填充作为初始提升，再反向提升径向同伦，并在边界指定原提升的反向过程。相对提升引理保证最终边界恰为原截面。这给出同一零值判据，而不要求整条纤维作为任意参数空间参与提升。

\begin{example}[连通度给出的自动延拓]\label{g2:exm:connectivity-extension-range}
	设 \(n\ge2\)，\(Y\) 为 \((n-1)\) 连通空间，\((X,A)\) 为 CW 对，\(f_A:A\to Y\) 已给定。选择新顶点的像，并用道路填入新一胞腔；随后 \(\pi_q(Y)=0\)（\(1\le q<n\)）使所有维数不超过 \(n\) 的新胞腔都可填入。因此 \(f_A\) 总可延到 \(X_n\)，首次可能出现的障碍位于 \((n+1)\) 胞腔，取值于 \(\pi_n(Y)\)。对具有 \((n-1)\) 连通非空纤维的纤维化，指定截面的相对延拓具有相同范围。
\end{example}

\subsection{基本群作用产生的系数系统}
\label{g2:subsec:1201:03}

单个胞腔只问一个球面类是否为零，所以任意基点运输都不会改变答案。若要把多个胞腔上的障碍相加，或比较两次选择造成的差，便必须说明这些类处于哪些群、不同群之间怎样识别。第5章的局部系统在这里重新出现，系数从纤维同调换成纤维同伦群。

\begin{definition}[简单空间]\label{g2:def:simple-space}
	道路连通空间 \(Y\) 若满足 \(\pi_1(Y)\) 交换，且它对每个 \(\pi_q(Y)\)（\(q\ge2\)）的基点运输作用都平凡，称为\term[jiandan-kongjian]{简单空间}[simple space]。对整数 \(q\ge1\)，若只要求这些作用在 \(2\le j\le q\) 时平凡，则称为 \(q\) 简单空间。\(q=1\) 时这个条件仅要求基本群交换。
\end{definition}

基本群对自身的作用是共轭，因此定义中的交换性恰好使一次群的基点运输也不受路径选择影响。由\cref{g2:prop:basepoint-action}，两条具有相同端点的道路给出的基点变换相差一个环路作用；对简单空间，这个差别消失。于是其各点的 \(\pi_q\) 可以作唯一的路径无关识别。每个单连通空间都是简单空间；\(S^1\) 也简单，因为其基本群为 \(\mathbb Z\)，高阶同伦群全为零。

\begin{proposition}[纤维同伦群的局部系统]\label{g2:prop:homotopy-local-system}
	设 \(p:E\to X\) 为基底道路连通、纤维非空且简单的 Hurewicz 纤维化。对每个 \(q\ge1\)，纤维同伦群组成阿贝尔群的局部系统 \(\Pi_q(p)\)：在 \(x\in X\) 处取群 \(\pi_q(F_x)\)，沿道路 \(\gamma:x\to y\) 的同态由纤维运输 \(T_\gamma:F_x\to F_y\) 诱导，并用纤维内部的路径识别基点。它只依赖于 \(\gamma\) 的定端点同伦类，且
	\begin{equation}\label{g2:eq:homotopy-local-transport}
		\Pi_q(p)(\gamma*\eta)=\Pi_q(p)(\eta)\Pi_q(p)(\gamma).
	\end{equation}
	若只处理次数至 \(q\)，纤维 \(q\) 简单已足够。
\end{proposition}
\begin{proof}
	第11章的道路提升把 \(T_\gamma\) 确定到自由同伦；反向道路给出同伦逆，定端点道路同伦也给出相应映射同伦。若为 \(F_x,F_y\) 各选一点，\(T_\gamma\) 的基点像未必等于已选终点基点，故还须沿 \(F_y\) 内的一条道路变换基点。简单性说明这样的基点变换与所选道路无关。改变 \(T_\gamma\) 的提升选择时，同伦的基点轨迹也只造成同一类已经消失的环路作用，因而群同态无歧义。

	纤维运输的拼接满足 \(T_{\gamma*\eta}\simeq T_\eta T_\gamma\)，诱导同态遂满足\eqref{g2:eq:homotopy-local-transport}；常值道路诱导恒等映射。这正是\cref{g2:def:local-system}的局部系统公理。\(q\ge2\) 时同伦群本来就交换，\(q=1\) 时交换性由简单性提供。对 Serre 纤维化，可通过第11章的 Hurewicz 替换及纤维弱同伦等价搬回这一系统，或直接对球面参数作上述运输。
\end{proof}

“纤维简单”消除的是同一纤维内的基点歧义，并不要求绕基底一周后的运输恒等。例如取反足映射 \(a:S^2\to S^2\)，其映射环面
\[
	E=(S^2\times I)/((z,1)\sim(a(z),0))\longrightarrow S^1
\]
具有单连通纤维，但基底正向绕行在 \(\pi_2(S^2)=\mathbb Z\) 上作用为乘 \(-1\)。这是因为 Hurewicz 同构把该作用对应到反足映射在二次同调上的度 \(-1\)。因此 \(\Pi_2(p)\) 是反号局部系统，不能以常值整数系数取代。定向 \(S^m\) 丛的情况则相反：其顶次运输保持定向，故 \(\Pi_m(p)\cong\mathbb Z\) 为常值系统；未定向时，它记录的正是纤维定向反转。

\ProseParagraphBreak

普通映射问题中，若 \(Y\) 简单，乘积纤维化 \(X\times Y\to X\) 的运输为恒等，因而系数就是常值群 \(\pi_q(Y)\)。若目标不简单，则已给的低维映射也会进入系数。例如一张定义在 \(X_2\) 上的映射决定
\[
	\rho:\pi_1(X)\longrightarrow\pi_1(Y),
\]
因为 \(X\) 的基本群由二骨架决定。固定 \(\rho\) 后，\(\pi_1(Y)\) 在 \(\pi_q(Y)\) 上的作用可沿 \(\rho\) 拉回，成为 \(X\) 上的局部系统。仅给出抽象群 \(\pi_q(Y)\)，还没有指定这组系数。下一节采用简单纤维而允许基底运输不平凡的条件，既保持阿贝尔上同调的运算，也包括非定向球丛这样的几何情形。

\ProseParagraphBreak

因此，第 \(q+1\) 层的障碍将由相对胞腔上链记录在
\[
	C^{q+1}(X,A;\Pi_q(p))
\]
中。每个胞腔的定向决定球面类的符号，道路运输决定不同胞腔的系数如何相加。下一节要证明这些值满足余圈条件，并确定修改已有延拓时产生的余边界；只有完成这两步，才能把逐胞腔判据提升为上同调中的不变量。映射与提升问题的阅读对照可参看\cite[第4.3节，Obstruction Theory]{Hatcher2002}及\cite[第18章，第5节，Obstruction theory]{May1999}。
% 纲要条目排为编号小节标题，并进入目录与书签；说明段仍保留于纲要注释中。

% !TeX root = ../../Book2.tex
% 纲要标题：### 12.2 障碍余圈及其不变量
% 纲要位置：几何学计划纲要.md，第 1478 行。
% 正文状态：初稿；尚未正式审核。

% 纲要内容开始
% 
% * 障碍上链的构造
% * 余圈条件的证明
% * 改变初始延拓时的余边界
% * 障碍类消失与延拓存在
% * 差异类与上同调可表性
% 
% 纲要内容结束

\section{障碍余圈及其不变量}
\label{g2:sec:1202}

一个边界球面类判断一个胞腔能否填入。对整个空间，胞腔的边界之间还有相互关系，而重新选择已经填好的部分也会改变下一步的障碍。上链记录逐胞腔的值，余圈条件记录这些值之间的关系，除去余边界则记录允许改变上一层后的不变量。本节逐一证明这三件事。

\subsection{障碍上链的构造}
\label{g2:subsec:1202:01}

固定道路连通 CW 对 \((X,A)\)、Hurewicz 纤维化 \(p:E\to X\)，以及 \(A\) 上的指定截面 \(s_A\)。假定纤维非空且简单，记上一节构造的局部系统为 \(\Pi_n(p)\)。对某个 \(n\ge1\)，设已经选定延拓 \(s:X_n=A\cup X^n\to E\)。

\ProseParagraphBreak

给每个相对 \((n+1)\) 胞腔选定向及特征映射 \(\chi_e:D^{n+1}\to X_{n+1}\)。沿圆盘的径向收缩把边界截面送到中心纤维，便得到一个球面映射 \(a_e:S^n\to F_{\chi_e(0)}\)。源球面取圆盘的正向边界；它的同伦类就是这个胞腔上的延拓障碍。

\begin{definition}[障碍上链]\label{g2:def:obstruction-cochain}
	在上述条件下，将每个有向相对 \((n+1)\) 胞腔 \(e\) 赋值为 \([a_e]\in\pi_n(F_{\chi_e(0)})\)，得到\term[zhangai-shanglian]{障碍上链}[obstruction cochain]
	\[
		c_{n+1}(s)\in C^{n+1}(X,A;\Pi_n(p)),\qquad c_{n+1}(s)(e)=[a_e].
	\]
	在 \(A\) 内的胞腔上取零。若通过道路把系数搬到选定的基点纤维，则同时按第5章的局部系数约定记录运输。
\end{definition}

这一赋值包含三项选择，但它们的作用可分别追踪。反转胞腔定向使球面类变号；改变圆盘上的径向提升，只产生纤维内同伦；改变连接基点的道路，则按 \(\Pi_n(p)\) 的运输变换系数。纤维简单性保证同一纤维内的基点道路不再产生歧义。因此得到的是一个确定的局部系数上链，而不是把不同纤维中的元素直接相加。

\ProseParagraphBreak

由\cref{g2:prop:section-cell-extension}，保持 \(s\) 在 \(X_n\) 上的全部值时，延拓到 \(X_{n+1}\) 的充要条件是
\begin{equation}\label{g2:eq:obstruction-cochain-vanishing}
	c_{n+1}(s)=0.
\end{equation}
这是上链逐项为零的条件。它与稍后出现的上同调类为零有不同含义。

\begin{example}[带一个关系的二维复形]\label{g2:exm:obstruction-relation-cochain}
	令 \(X=S^1\cup_m e^2\)，二胞腔沿圆周的 \(m\) 次绕行附着，\(m\ge2\)。给定一次骨架映射 \(f:S^1\to S^1\)，其次数为 \(k\)。把它看成乘积纤维化的截面，则
	\[
		c_2(f)(e^2)=mk\in\pi_1(S^1)=\mathbb Z.
	\]
	因为 \(m\ne0\)，保持这个 \(f\) 延拓的条件是 \(k=0\)。另一方面，胞腔上链的余边界 \(C^1(X;\mathbb Z)\to C^2(X;\mathbb Z)\) 是乘 \(m\)，所以 \([c_2(f)]=0\) 对所有 \(k\) 都成立。允许改变一胞腔上的映射后，确实总能把它换为常值再延拓。这个例子先把两种零值判据的差别显示出来。
\end{example}

\subsection{余圈条件的证明}
\label{g2:subsec:1202:02}

虽然 \(s\) 尚未定义在 \((n+1)\) 胞腔内部，障碍值仍须满足 \((n+2)\) 胞腔给出的关系。关键是：整个特征圆盘可缩，它的拉回纤维化有截面；这个截面未必与已有边界数据相同，但足以证明障碍值在胞腔边界上的总和为零。

\begin{lemma}[障碍余圈]\label{g2:lem:obstruction-cocycle}
	设 \((X,A)\) 是道路连通 CW 对，\(p:E\to X\) 为纤维非空且简单的 Hurewicz 纤维化。对 \(n\ge1\) 及截面 \(s:X_n\to E\)，有
	\[
		\delta c_{n+1}(s)=0\quad\text{于 }C^{n+2}(X,A;\Pi_n(p)).
	\]
\end{lemma}
\begin{proof}
	先把逐胞腔的定义扩大到相对圆盘。给定带基点映射 \(\Phi:(D^{n+1},S^n)\to(X_{n+1},X_n)\)，用 \(s\Phi\) 作为边界提升，沿 \(\Phi\) 给出的径向同伦提升到一个纤维，记所得类为 \(\omega_s[\Phi]\)。相对同伦的提升说明它只依赖相对同伦类；将两个圆盘沿加法方向拼接，再提升同一拼接，说明 \(\omega_s\) 是群同态。若 \(\Phi\) 本身已有一个与指定边界相同的提升，则所得球面可填，\(\omega_s[\Phi]=0\)。

	为了同时处理运输，拉回到 \(X\) 的万有覆盖 \(\widetilde X\)。局部系统在那里成为常值系统；对每个骨架记其全部逆像为 \(\widetilde X_r\)。当 \(r\ge2\) 时，\(X_r\hookrightarrow X\) 在基本群上为同构，所以 \(\widetilde X_r\) 单连通。对 \(n\ge2\)，空间对 \((\widetilde X_{n+1},\widetilde X_n)\) 只附着 \((n+1)\) 胞腔，故为 \(n\) 连通；\cref{g2:thm:hurewicz-relative}给出
	\[
		\pi_{n+1}(\widetilde X_{n+1},\widetilde X_n)\xrightarrow{\cong}
		H_{n+1}(\widetilde X_{n+1},\widetilde X_n).
	\]
	特征圆盘在这张同构下成为相应的有向胞腔。因此 \(\omega_s\) 在相对同调上的表示，恰好是障碍上链的求值，而不是另选的一张同态。

	现在取一个 \((n+2)\) 胞腔的特征映射 \(g:D^{n+2}\to\widetilde X_{n+2}\)，并在边界上选基点。其胞腔边界由以下路径得到：先取 \(g|_{S^{n+1}}\) 的球面类，再把它送入 \(\pi_{n+1}(\widetilde X_{n+1},\widetilde X_n)\)，最后用上述 Hurewicz 同构。这与同调边界一致，因为 Hurewicz 映射与边界及包含相容，见\cref{g2:prop:hurewicz-naturality}。

	把纤维化拉回 \(D^{n+2}\)，可选一个在边界基点处等于 \(s\) 的截面。它给出 \(g|_{S^{n+1}}\) 的整体提升。将这个球面写成 \(D^{n+1}/S^n\) 后，所得相对圆盘在其边界恒取选定基点，提升在该边界也恒等于 \(s\) 的值。因此这一相对类被 \(\omega_s\) 送到零。换言之，\(c_{n+1}(s)\) 在任一 \((n+2)\) 胞腔的胞腔边界上取零。

	\(n=1\) 时，\(\widetilde X_1\) 未必单连通，需要把上述相对同伦到胞腔链的步骤直接说明。对二维相对圆盘，将它经过的二维胞腔表示为有向小圆盘，用细弧连接到边界基点；剩余部分映到 \(\widetilde X_1\)。沿细弧切开后，每个小圆盘贡献其胞腔障碍的正值或负值，剩余部分已有截面，贡献为零。重新粘合时，边上的道路成对反向出现；改变连接弧只把相应纤维环路共轭。由于 \(\pi_1(F)\) 交换，共轭消失，所得总值只依赖各二维胞腔的有向次数。这就证明 \(\omega_s\) 经由 \(H_2(\widetilde X_2,\widetilde X_1)\) 分解。上述小圆盘表示来自相对胞腔逼近及圆盘映射的次数分解，与第3章计算胞腔关联数所用的分解相同。

	对三胞腔的附着球面使用这一分解，其每条内部边出现两次、方向相反；也可沿前一段的圆盘提升，直接看出相应相对类的 \(\omega_s\) 为零。因此 \(n=1\) 同样得到 \(\delta c_2(s)=0\)。最后把覆盖中的计算降回 \(X\)。改变提升胞腔会同时改变其中心纤维的识别，恰好插入 \(\Pi_n(p)\) 的运输；第5章的局部系数胞腔余边界正是这一等变求值。因此结论对原来的局部系数成立。\(A\) 内的胞腔原已有截面，所有相应值为零，故上述论证也保留相对条件。
\end{proof}

由引理可定义
\begin{equation}\label{g2:eq:obstruction-class}
	o_{n+1}(s)=[c_{n+1}(s)]\in H^{n+1}(X,A;\Pi_n(p)).
\end{equation}
这称为当前截面的\term[zhangai-lei]{障碍类}[obstruction class]。上同调群描述障碍可能出现的位置；\(o_{n+1}(s)\) 是当前数据给出的具体元素。一个非零群完全可能包含等于零的障碍类。

\subsection{改变初始延拓时的余边界}
\label{g2:subsec:1202:03}

设 \(s,t\) 都定义在 \(X_n\)，并在 \(X_{n-1}\) 上相同。对一个 \(n\) 胞腔，将 \(t\) 的圆盘取正向、\(s\) 的圆盘取负向，沿共同边界粘成 \(S^n\)，再利用圆盘上的纤维同伦模型送入中心纤维。所得类记为
\[
	D_n(s,t)(e^n)\in\Pi_n(p)_{\chi_e(0)}.
\]
这些值组成\term[chayi-shanglian]{差异上链}[difference cochain] \(D_n(s,t)\in C^n(X,A;\Pi_n(p))\)。它记录把 \(s\) 改成 \(t\) 时，在每个 \(n\) 胞腔中实际加入的球面类。

\ProseParagraphBreak

若 \(s,t\) 只在 \(X_{n-1}\) 上由指定截面同伦 \(h\) 相连，则用 \(h\) 填侧面，将上下两张圆盘联成柱体的整个边界，同样得到 \(D_n(s,t;h)\)。为固定符号，柱体取先胞腔、后区间的乘积定向。若 \(H\) 是两端用 \(s,t\)、侧面用 \(h\) 的部分截面，则
\begin{equation}\label{g2:eq:difference-prism-sign}
	D_n(s,t;h)(e^n)=(-1)^n c_{n+1}(H)(e^n\times I).
\end{equation}
因为 \(\partial(e^n\times I)=\partial e^n\times I+(-1)^n(e^n\times\{1\}-e^n\times\{0\})\)，这一约定恰好使上端为 \(t\)、下端为 \(-s\)。

\begin{proposition}[差异与余边界]\label{g2:prop:obstruction-difference}
	设 \((X,A)\) 为道路连通 CW 对，\(p:E\to X\) 为纤维非空且简单的 Hurewicz 纤维化，\(n\ge1\)。若 \(s,t:X_n\to E\) 是在 \(A\) 上相同的截面，\(h\) 是它们在 \(X_{n-1}\) 上固定 \(A\) 的截面同伦，则
	\begin{equation}\label{g2:eq:obstruction-difference}
		\delta D_n(s,t;h)=c_{n+1}(t)-c_{n+1}(s).
	\end{equation}
	此外，给定 \(s\) 和任意 \(b\in C^n(X,A;\Pi_n(p))\)，都能选到在 \(X_{n-1}\) 上与 \(s\) 相同的 \(t\)，使 \(D_n(s,t)=b\)。
\end{proposition}
\begin{proof}
	令 \(W=X_{n+1}\)，在 \(W\times I\) 上拉回纤维化。部分截面 \(H\) 已在 \(X_n\times\partial I\)、\(X_{n-1}\times I\) 和 \(A\times I\) 上定义，这包含相对柱体的 \(n\) 骨架。由上一引理，其 \((n+1)\) 次障碍在每个 \((n+2)\) 胞腔 \(e^{n+1}\times I\) 的边界上取零。因此
	\[
		0=\sum_{e^n}[e^{n+1}:e^n]c_{n+1}(H)(e^n\times I)
		+(-1)^{n+1}\bigl(c_{n+1}(t)-c_{n+1}(s)\bigr)(e^{n+1}).
	\]
	有局部系数时，求和中同时沿各附着道路运输；两端的系数用区间方向的恒等运输识别。代入式\eqref{g2:eq:difference-prism-sign}并乘 \((-1)^n\)，便得到所需公式。

	为实现指定的 \(b\)，在每个相对 \(n\) 胞腔内部取一个小圆盘。在纤维模型中，把该圆盘的边界压成一点，插入代表 \(b(e^n)\) 的球面映射；原来通向小圆盘中心的值通过一条纤维道路与球面基点连接。等价地，这是保持外边界的挤压映射 \(D^n\to D^n\vee S^n\)，在第一因子沿用原映射，在第二因子加入所选球面类。简单性使连接道路不影响所得类，方向取为 \(t-s\)，于是差异正是 \(b(e^n)\)。对 \(n=1\)，此操作就是在原道路中插入指定环路。

	纤维同伦模型转换回原丛后，若外边界只相等到纤维同伦，按\cref{g2:prop:section-cell-extension}证明中的相对提升把它修正为实际相等。各胞腔内部的修改相互独立，在 \(X_{n-1}\) 及 \(A\) 上全不改变，故由 CW 弱拓扑粘合成 \(t\)。它满足 \(D_n(s,t)=b\)。
\end{proof}

因此，保持 \(X_{n-1}\) 上的数据时，\(c_{n+1}(s)\) 恰好可以改变一个任意余边界。障碍类正是这组修改下保持不变的量。普通映射的对应叙述见\cite[第18章第5节，Obstruction theory]{May1999}；本节的局部运输与柱体公式给出了截面版本的同一机制。

\subsection{障碍类消失与延拓存在}
\label{g2:subsec:1202:04}

现在可以把“修改后填入”的几何要求写成一个上同调判据。

\begin{theorem}[逐层延拓判据]\label{g2:thm:obstruction-extension}
	设 \((X,A)\) 为道路连通 CW 对，\(p:E\to X\) 为纤维非空且简单的 Hurewicz 纤维化，\(n\ge1\)，且 \(s:X_n\to E\) 为截面。则以下两条件等价：
	\begin{enumerate}
		\item \(o_{n+1}(s)=0\in H^{n+1}(X,A;\Pi_n(p))\)；
		\item \(s|_{X_{n-1}}\) 可保持不变地延拓为 \(X_{n+1}\) 上的截面。
	\end{enumerate}
	若要求保持原来的整个 \(s|_{X_n}\)，充要条件是上链 \(c_{n+1}(s)=0\)。
\end{theorem}
\begin{proof}
	若 \(o_{n+1}(s)=0\)，选 \(b\in C^n(X,A;\Pi_n(p))\) 使 \(\delta b=-c_{n+1}(s)\)。差异命题给出与 \(s\) 在 \(X_{n-1}\) 上相同的 \(t\)，满足
	\[
		c_{n+1}(t)=c_{n+1}(s)+\delta b=0.
	\]
	单胞腔判据把 \(t\) 延拓到 \(X_{n+1}\)。反过来，若有这样的延拓，记它在 \(X_n\) 上的限制为 \(t\)，则 \(c_{n+1}(t)=0\)，而差异公式使 \(c_{n+1}(s)=-\delta D_n(s,t)\)，故障碍类为零。最后一句直接是\cref{g2:eq:obstruction-cochain-vanishing}。
\end{proof}

\begin{theorem}[首障碍]\label{g2:thm:primary-obstruction}
	设 \((X,A)\) 为道路连通 CW 对，\(p:E\to X\) 为纤维非空、简单且 \((n-1)\) 连通的 Hurewicz 纤维化，\(n\ge1\)，并给定截面 \(s_A\)。则 \(s_A\) 总可延到 \(X_n\)，且
	\[
		o_{n+1}(p,s_A)\in H^{n+1}(X,A;\Pi_n(p))
	\]
	与所选 \(X_n\) 上的延拓无关。它为零当且仅当 \(s_A\) 可延到 \(X_{n+1}\)。这个\term[shouzhangai]{首障碍}[primary obstruction]对保持给定截面的 CW 对拉回自然。若相对胞腔维数均不超过 \(n+1\)，它为零便等价于全空间上的截面存在。
\end{theorem}
\begin{proof}
	非空性允许选择新顶点的值，道路连通性允许填入新一胞腔；\(\pi_j(F)=0\)（\(0<j<n\)）再使维数至 \(n\) 的胞腔都可以填入。故存在所需初始延拓。

	取两张这样的截面 \(s,t\)。在 \(X_{n-1}\) 上构造固定 \(A\) 的截面同伦，只需处理 \(k\) 胞腔乘区间，\(k\le n-1\)；其边界问题取值于 \(\pi_k(F)\)。当 \(k=0\) 时由道路连通性填入，其余均由低阶同伦群消失填入。于是得到前一命题所需的 \(h\)，差异公式证明两张障碍上圈同属一个上同调类。

	若这个类为零，逐层延拓判据给出延到 \(X_{n+1}\) 的截面；反之，取实际延拓的限制来计算，障碍上链本身为零。对胞腔映射的拉回，在每个相对胞腔上，边界球面复合及运输与原构造相容；将复合映射分解成有向胞腔贡献，就得到局部系数上链上的拉回公式。一般映射先相对胞腔逼近，拉回纤维化在基底同伦下纤维同伦等价，故得到相同的障碍类。这证明自然性。最后，在所述维数范围内 \(X_{n+1}=X\)，所以无需再处理更高胞腔。
\end{proof}

纤维较低同伦群的消失在这里承担两个不同任务：它使初始截面存在，也使更低骨架上的初始选择同伦唯一。缺少第二项时，第 \(n+1\) 次障碍类仍对固定的低层数据有意义，却可能随着更低层的选择改变。

\begin{corollary}[由障碍群消失得到延拓]\label{g2:cor:obstruction-groups-vanish}
	设 \((X,A)\) 为道路连通 CW 对，\(p:E\to X\) 为纤维非空且简单的 Hurewicz 纤维化。若
	\[
		H^{q+1}(X,A;\Pi_q(p))=0\qquad(q\ge1),
	\]
	则每个 \(A\) 上的截面都可延拓到 \(X\)。对有限维相对 CW 对，只需核对实际出现的次数。
\end{corollary}
\begin{proof}
	先延拓到 \(X_1\)。在第 \(q\) 步，障碍群为零使当前障碍类为零，所以可以保持 \(X_{q-1}\) 不变，修改 \(q\) 胞腔上的值，再延拓到 \(X_{q+1}\)。一个固定的 \(k\) 胞腔只会在第 \(k\) 步最后修改一次，以后所有步骤都保留它。取这些最终值，便在每个特征圆盘上得到连续且边界相容的截面；CW 弱拓扑使它们定义整个 \(X\) 上的连续截面。
\end{proof}

\subsection{差异类与上同调可表性}
\label{g2:subsec:1202:05}

截面存在以后，还可以问它是否同伦唯一。这一次要填的是两张截面之间的柱体，因此障碍所在的上同调次数降低一。

\begin{proposition}[同伦的差异类]\label{g2:prop:section-difference-class}
	设 \((X,A)\)、\(p:E\to X\) 满足上一推论的空间与纤维条件，\(s,t\) 为在 \(A\) 上相同的全局截面。给定它们在 \(X_{n-1}\) 上固定 \(A\) 的截面同伦 \(h\)，\(n\ge1\)，则 \(D_n(s,t;h)\) 为余圈。其类
	\[
		\Delta_n(s,t;h)\in H^n(X,A;\Pi_n(p))
	\]
	为零，当且仅当 \(h|_{X_{n-2}}\) 可延拓为 \(s|_{X_n}\) 与 \(t|_{X_n}\) 之间的截面同伦。\(n=1\) 时 \(X_{-1}=A\)，所以只固定 \(A\) 上的同伦。
\end{proposition}
\begin{proof}
	两张全局截面的障碍上链均为零，所以\cref{g2:eq:obstruction-difference}给出 \(\delta D_n=0\)。在空间对
	\[
		\bigl(X\times I,\ X\times\partial I\cup A\times I\bigr)
	\]
	上考虑给定的两端与部分同伦。相对胞腔为 \(e^k\times(0,1)\)，维数为 \(k+1\)。相对链的边界去掉两端后只有 \(\partial e^k\times I\)，所以其局部系数上同调自然识别为 \(H^{j-1}(X,A;\Pi_n(p))\)。这就是柱体相对方向使次数降低一的原因。

	用式\eqref{g2:eq:difference-prism-sign}识别同伦障碍，逐层延拓判据便说明：其类为零，恰好允许修改内部的 \(n\) 维胞腔，即 \(e^{n-1}\times I\) 上的同伦，同时固定更低柱体与全部两端，再填入 \(e^n\times I\)。这正是结论所述的延拓，而不是要求保留原来的整个 \(h|_{X_{n-1}}\)。
\end{proof}

如果纤维 \((n-1)\) 连通，上述 \(h\) 在所需低次中总能找到。其不同选择在更低的柱体上又可同伦，因为相关同伦群仍在消失范围内；差异公式于是使 \(\Delta_n\) 的上同调类与这些低次选择无关。接下来两种情形中，这个差异类已经记录全部自由度。

\begin{theorem}[截面同伦类的参数]\label{g2:thm:section-torsor}
	设 \((X,A)\) 为道路连通 CW 对，\(p:E\to X\) 为具有非空、简单、\((n-1)\) 连通纤维的 Hurewicz 纤维化，\(n\ge1\)，并假设指定截面 \(s_A\) 至少有一个全局延拓。以下任一条件成立时，延拓的相对截面同伦类形成 \(H^n(X,A;\Pi_n(p))\) 上的主齐性集：
	\begin{enumerate}
		\item 相对胞腔维数均不超过 \(n\)；
		\item 纤维的全部正阶同伦群除 \(\pi_n\) 外均为零。
	\end{enumerate}
	选定一个延拓 \(s_0\) 后，\([t]\mapsto\Delta_n(s_0,t)\) 给出与该上同调群的双射。
\end{theorem}
\begin{proof}
	低阶同伦群消失使各截面在 \(X_{n-1}\) 上可由相对同伦相连。选定这样的同伦后，逐胞腔球面的拼接给出
	\[
		\Delta_n(s,t)+\Delta_n(t,u)=\Delta_n(s,u),\qquad
		\Delta_n(s,s)=0.
	\]
	改变辅助低次同伦只加一个余边界：将两次同伦再放到一个参数区间上比较，在维数低于 \(n\) 的地方由纤维连通度填充，剩余柱体边界正是差异公式。因而这些等式确实是上同调中的等式。

	若 \(\Delta_n(s,t)=0\)，上一命题给出 \(X_n\) 上的相对同伦。条件 I 下已经完成全空间；条件 II 下，继续填入更高的 \(e^k\times I\) 所遇系数 \(\pi_k(F)\) 均为零，所以同伦延拓到全部 \(X\)。反过来，全局同伦当然使每一级差异障碍为零。这证明所列映射单射。

	给定任意上同调类，取余圈 \(b\in C^n(X,A;\Pi_n(p))\) 表示它。在 \(s_0\) 的各 \(n\) 胞腔中插入 \(b\) 的球面代表，得到 \(t_n\)，并有 \(D_n(s_0,t_n)=b\)。条件 I 下它已经是全局截面；条件 II 下，由 \(\delta b=0\) 及 \(c_{n+1}(s_0)=0\)，得到 \(c_{n+1}(t_n)=0\)，先延到 \(X_{n+1}\)，再利用更高同伦群消失继续延拓。所得截面给出指定的差异类，证明满射。

	因此上同调群通过加入差异类自由且传递地作用于延拓同伦类。这样的集合称为\term[zhuqixingji]{主齐性集}[torsor]：任意两个元素之间有唯一的群元素作为差，但在选择 \(s_0\) 前没有一个自然的零元。
\end{proof}

\begin{theorem}[上同调的可表性]\label{g2:thm:em-representability}
	设 \(G\) 是阿贝尔群，\(n\ge1\)，\((X,A)\) 为 CW 对。取带基点 CW 模型 \(K(G,n)\)，并令 \(\iota_n\in H^n(K(G,n);G)\) 为在 \(\pi_n\cong G\) 上对应恒等映射的类。则有自然双射
	\[
		[(X,A),(K(G,n),*)]\xrightarrow{\cong}H^n(X,A;G),\qquad [f]\longmapsto f^*\iota_n.
	\]
	左端是把 \(A\) 映到基点的映射在固定 \(A\) 的同伦下的类；无相对条件时得到 \([X,K(G,n)]\cong H^n(X;G)\)。
\end{theorem}
\begin{proof}
	第9章已经构造了 \(K(G,n)\)。\(n\ge2\) 时，Hurewicz 定理给出 \(H_n(K;\mathbb Z)\cong G\) 以及低正次同调消失，故万有系数定理把 \(H^n(K;G)\) 识别为 \(\operatorname{Hom}(G,G)\)。\(n=1\) 时，用一次 Hurewicz 定理及 \(G\) 交换得到相同识别；\(H_0=\mathbb Z\) 不产生 Ext 项。因此 \(\iota_n\) 定义明确。

	对乘积纤维化 \(X\times K(G,n)\to X\)，常值映射是指定边界值的一个延拓。纤维为 \((n-1)\) 连通简单空间，且只剩一个正阶同伦群，所以上一定理的条件 II 给出由 \(H^n(X,A;G)\) 参数化的全部映射同伦类。\(X\) 不连通时逐分支应用即可。

	还要验证所得参数确实是 \(f^*\iota_n\)。先把 \(f\) 在 \(X_{n-1}\) 上同伦为常值，并由同伦延拓性质将这段同伦延到全空间。对每个相对 \(n\) 胞腔，此时 \(f\) 压去边界后给出 \(S^n\to K(G,n)\)。差异上链在该胞腔上的值就是此球面代表的 \(G\) 中元素；而 \(\iota_n\) 对这个球面基本类的求值也正是该元素，因为它对应 \(G\) 上的恒等映射。因此两个胞腔上圈相同，给出所列拉回公式。公式与空间对映射的复合相容，所以双射自然。
\end{proof}

这一结论把一个抽象上同调类变成一张实际映射。下一节正要利用它：把一个阶段的粘合类表示为到 \(K(G,n+1)\) 的映射，再取该映射的同伦纤维，便能把相应的 \(G\) 放回所需次数。关于可表性、Postnikov 阶段与障碍问题之间的关系，可参看\cite[第4.3节，Connections with Cohomology]{Hatcher2002}。
% 纲要条目排为编号小节标题，并进入目录与书签；说明段仍保留于纲要注释中。

% !TeX root = ../../Book2.tex
% 纲要标题：### 12.3 Postnikov 塔的构造
% 纲要位置：几何学计划纲要.md，第 1486 行。
% 正文状态：初稿；尚未正式审核。

% 纲要内容开始
% 
% * 逐次消去高阶同伦群
% * Eilenberg–MacLane 纤维
% * k-不变量和阶段之间的粘合
% * 同伦群相同的不同三型
% 
% 纲要内容结束

\section{Postnikov 塔的构造}
\label{g2:sec:1203}

前两节把映射的构造分成逐个胞腔的延拓。本节从另一端分解问题：保留目标空间的低阶同伦群，逐次舍去高阶同伦群。第11章已用这样的截段证明有限生成性；现在要说明相邻截段怎样接在一起，以及这部分接合信息为何不能从同伦群的清单中读出。

\subsection{逐次消去高阶同伦群}
\label{g2:subsec:1203:01}

设 \(X\) 为单连通带基点 CW 复形。\cref{g2:lem:finite-homotopy-truncations}构造了 \(X\hookrightarrow X[n]\)：依次附着 \((n+2),(n+3),\ldots\) 维胞腔，消去 \(\pi_{n+1},\pi_{n+2},\ldots\)，同时保留 \(\pi_i\)（\(i\le n\)）。注意这与取 \(n\) 骨架的方向相反。取骨架是删去高维胞腔，可能重新产生高阶同伦；这里则添加高维胞腔，为原来不能填充的球面提供填充。

\begin{definition}[Postnikov 塔]\label{g2:def:postnikov-tower}
	设 \(X\) 为道路连通带基点空间。一列带基点映射
	\[
		\cdots\xrightarrow{p_{n+1}}P_nX\xrightarrow{p_n}P_{n-1}X\longrightarrow\cdots\longrightarrow P_1X
	\]
	连同相容映射 \(a_n:X\to P_nX\)，若满足 \(p_na_n=a_{n-1}\)、\((a_n)_*:\pi_i(X)\xrightarrow{\cong}\pi_i(P_nX)\)（\(i\le n\)），以及 \(\pi_i(P_nX)=0\)（\(i>n\)），就称为 \(X\) 的\term[Postnikov ta]{Postnikov 塔}[Postnikov tower]。空间 \(P_nX\) 称为第 \(n\) 阶段；高于 \(n\) 的同伦群全为零的空间也称为一个\term[n xing]{\(n\)-型}[\(n\)-type]。
\end{definition}

Postnikov 塔\footnote{Mikhail Mikhailovich Postnikov（1927—2004），苏联、俄罗斯数学家，主要研究代数拓扑与同伦理论，长期任教于莫斯科大学。}的下标记录保留的最高同伦次数。对单连通空间，可取 \(P_1X=\{*\}\)。第一个有可能非平凡的阶段是 \(P_2X\)，它只保留 \(\pi_2(X)\)。\symidx{\(P_nX\)：空间的第 \(n\) 个 Postnikov 阶段}

\begin{theorem}[Postnikov 塔的存在]\label{g2:thm:postnikov-tower-existence}
	每个单连通带基点 CW 复形 \(X\) 都有 Postnikov 塔，且可使每张 \(p_n\)（\(n\ge2\)）为 Hurewicz 纤维化，各阶段具有 CW 同伦型。
\end{theorem}
\begin{proof}
	先用 \(X[n]\) 构造相容的 CW 阶段。包含 \(X\to X[n-1]\) 可延拓为 \(q_n:X[n]\to X[n-1]\)：在一个新加的 \(d\) 胞腔上，\(d\ge n+2\)，延拓障碍属于 \(\pi_{d-1}(X[n-1])=0\)。逐胞腔作延拓，并用相对胞腔逼近保持 \(X\) 上的原映射，就得到 \(q_n|_X=\operatorname{id}_X\) 后接包含。各 \(X[n]\) 已有定义所需的同伦群。

	再由下往上作路径纤维化替换。取 \(P_1X=\{*\}\)，并令 \(j_1:X[1]\to P_1X\) 为同伦等价。假定已有同伦等价 \(j_{n-1}:X[n-1]\to P_{n-1}X\)。把复合 \(f_n=j_{n-1}q_n\) 替换为
	\[
		P_nX=\{(z,\gamma):z\in X[n],\ \gamma:I\to P_{n-1}X,\ \gamma(0)=f_n(z)\},
		\qquad p_n(z,\gamma)=\gamma(1).
	\]
	\cref{g2:prop:fibration-replacement}说明 \(p_n\) 是 Hurewicz 纤维化，而 \(j_n(z)=(z,\text{常值道路 }f_n(z))\) 是同伦等价。令 \(a_n=j_n|_X\)，便有严格等式
	\[
		p_na_n=j_{n-1}q_n|_X=j_{n-1}|_X=a_{n-1}.
	\]
	所以替换没有破坏塔的相容性，也没有改变任何阶段的同伦群。
\end{proof}

这项构造有许多选择，但阶段作为 \(X\) 的截段具有确定的同伦意义。若 \(Y\) 是带基点 \(n\)-型，则限制给出双射
\begin{equation}\label{g2:eq:postnikov-truncation-universal}
	[X[n],Y]_*\xrightarrow{\cong}[X,Y]_*.
\end{equation}
存在性来自新增胞腔的边界映入 \(Y\) 后总能填充，因为边界维数至少为 \(n+1\)。唯一性则把同一延拓问题用于圆柱：一个相对 \(d\) 胞腔对应的同伦延拓障碍在 \(\pi_d(Y)\) 中，仍为零。无限多个胞腔的延拓按弱拓扑粘合，得到全局映射与同伦。因此两个截段之间可延拓 \(X\) 上的相同映射，得到互逆到同伦的比较；这也说明将 \(X[n]\) 换成同伦等价的 \(P_nX\) 只改变模型。

\ProseParagraphBreak

特别地，\(P_2X\) 是 \(K(\pi_2(X),2)\) 的同伦型；如果 \(X\) 已经是 \(N\)-型，则 \(a_N:X\to P_NX\) 在全部同伦群上为同构，CW 模型之间的 Whitehead 定理便使它成为同伦等价。阶段保留的是整个低阶同伦结构，而不仅是某几个群的抽象副本。

\subsection{Eilenberg–MacLane 纤维}
\label{g2:subsec:1203:02}

相邻阶段只相差一个同伦群，纤维正好把这个差别单独装出来。

\begin{proposition}[阶段之间的纤维]\label{g2:prop:postnikov-stage-fiber}
	设 \(X\) 为单连通带基点 CW 复形，取上一构造的纤维化塔。对 \(n\ge2\)，\(p_n:P_nX\to P_{n-1}X\) 的纤维 \(F_n\) 道路连通，并满足
	\[
		\pi_i(F_n)=0\quad(i\ne n),\qquad \pi_n(F_n)\cong\pi_n(X).
	\]
	因此它的 CW 逼近为 \(K(\pi_n(X),n)\)，阶段纤维序列可记为
	\begin{equation}\label{g2:eq:postnikov-stage-fibration}
		K(\pi_n(X),n)\longrightarrow P_nX\xrightarrow{p_n}P_{n-1}X.
	\end{equation}
\end{proposition}
\begin{proof}
	总空间与基底均单连通，长正合列的末段先给出纤维道路连通。对 \(i<n\)，映射 \((p_n)_*\) 是由两张 \(a_n\) 给出的同构；在 \(i=n\) 时，基底的 \(\pi_n\) 与 \(\pi_{n+1}\) 均为零。因此长正合列给出
	\[
		\pi_i(F_n)=0\ (i<n),\qquad \pi_n(F_n)\xrightarrow{\cong}\pi_n(P_nX)=\pi_n(X).
	\]
	对 \(i>n\)，两阶段对应的同伦群全为零，纤维群也为零。CW 逼近保留这些群；第9章的 Eilenberg–MacLane 定义及第11章对其模型唯一性的证明给出最后的识别。
\end{proof}

式\eqref{g2:eq:postnikov-stage-fibration}并没有说 \(P_nX\) 是基底与纤维的乘积。此前的 Hopf 丛已经表明，同样的基底与纤维可能产生不同的总空间。这里纤维只剩一个同伦群，使粘合信息能够由一个上同调类记录，但仍需要把这个类找出来。

\ProseParagraphBreak

整座纤维化塔还保留了 \(X\) 的弱同伦型。说明这一点有助于理解为什么逐阶研究最终足以恢复全部同伦信息。

\begin{proposition}[从塔恢复弱同伦型]\label{g2:prop:postnikov-limit}
	对\cref{g2:thm:postnikov-tower-existence}中的纤维化塔，逆极限
	\[
		P_\infty X=\{(z_1,z_2,\ldots)\in\textstyle\prod_nP_nX:p_n(z_n)=z_{n-1}\}
	\]
	带子空间拓扑。相容映射给出 \(a:X\to P_\infty X\)，它是弱同伦等价。
\end{proposition}
\begin{proof}
	固定 \(i\ge1\)。同伦群系统 \(\pi_i(P_nX)\) 从 \(n\ge i\) 起稳定为 \(\pi_i(X)\)。先证明
	\[
		\pi_i(P_\infty X)\longrightarrow\varprojlim_n\pi_i(P_nX)
	\]
	满射。给定右边相容的同伦类，逐层选带基点代表 \(f_n:S^i\to P_nX\)。已有 \(f_{n-1}\) 后，\(p_nf_n\) 与 \(f_{n-1}\) 带基点同伦；相对基点提升这段同伦，便能调整 \(f_n\) 使 \(p_nf_n=f_{n-1}\) 严格成立。所有坐标合在一起给出到逆极限的连续映射。

	再设 \(f:S^i\to P_\infty X\) 在每个坐标上都零伦。取足够大的 \(N\)，使 \(\pi_{i+1}(P_nX)\to\pi_{i+1}(P_{n-1}X)\) 在 \(n>N\) 时全为同构。先选 \(f_N\) 的零伦 \(G_N:D^{i+1}\to P_NX\)。假定已选 \(G_{n-1}\)，任选 \(f_n\) 的零伦 \(G_n\)。映射 \(p_nG_n\) 与 \(G_{n-1}\) 在边界相同，将它们沿边界反向粘合，得到 \(P_{n-1}X\) 中的一个 \((i+1)\) 球面类。上述同伦群满射允许从 \(P_nX\) 中提升这个类；在 \(G_n\) 的圆盘内部接入相反的球面代表，即可使两张圆盘映射相对边界同伦。相对边界的同伦提升又将其改成严格等式 \(p_nG_n=G_{n-1}\)。递推所得相容零伦给出 \(f\) 在逆极限中的零伦。前 \(N-1\) 个坐标由投影取得，所以舍去有限个初始阶段不会遗漏条件。

	同样的道路构造处理道路分支：各阶段道路连通，而 \(\pi_1(P_nX)=0\) 使逐层选择的道路可以调整为相容道路。因此逆极限道路连通。综上，\(\pi_i(P_\infty X)\) 与稳定群 \(\pi_i(X)\) 同构，且这些同构由 \(a\) 给出。
\end{proof}

逆极限本身未必是 CW 复形，所以这里结论写为弱同伦等价。取它的 CW 逼近后，Whitehead 定理便恢复 \(X\) 的 CW 同伦型。上面的证明实际使用了各次同伦群最终稳定；一般逆极限不能直接把同伦群逐项取极限。

\subsection{k-不变量和阶段之间的粘合}
\label{g2:subsec:1203:03}

前节的可表性定理把 \(H^{n+1}(B;G)\) 中的类表示为映射 \(B\to K(G,n+1)\)。给出这样的映射 \(\kappa\)，就可以拉回路径纤维化，构造
\begin{equation}\label{g2:eq:k-invariant-path-model}
	E_\kappa=\{(b,\lambda):b\in B,\ \lambda(0)=\kappa(b),\ \lambda(1)=*\}
	\longrightarrow B.
\end{equation}
纤维是 \(\Omega K(G,n+1)\)，其唯一非零同伦群是 \(\pi_n\cong G\)。道路从 \(\kappa(b)\) 指向基点，这正是第8章规定的 \(\operatorname{hofib}(\kappa)\) 模型。它拉回的是终点固定的路径纤维化 \(P^{\mathrm{in}}K\xrightarrow{\operatorname{ev}_0}K\)。

\ProseParagraphBreak

关键问题是：每张阶段纤维化是否都能由这种形式得到？下面用一个实际构造证明，而后再把构造得到的映射识别为障碍类。

\begin{lemma}[单一同伦纤维的粘合映射]\label{g2:lem:principal-em-fibration}
	设 \(q:Z\to B\) 是单连通带基点 CW 复形之间的映射，\(n\ge2\)，其同伦纤维具有 \(K(G,n)\) 的弱同伦型，其中 \(G\) 为阿贝尔群。则存在带基点映射 \(\kappa:B\to K(G,n+1)\) 和覆盖 \(B\) 的弱同伦等价
	\[
		j:Z\longrightarrow E_\kappa,\qquad \operatorname{pr}_B j=q.
	\]
	可以按指定的 \(\pi_n(\operatorname{hofib}q)\cong G\) 选择模型，使纤维比较在这个群上为恒等。
\end{lemma}
\begin{proof}
	先假设 \(q\) 胞腔。一般情形取带基点胞腔逼近 \(q'\) 及同伦 \(H:q\simeq q'\)；为 \(q'\) 构造比较后，在所得道路前接上 \(\kappa H(z,-)\)，并将基底坐标改为 \(q(z)\)，便得到原映射的比较。令 \(M=M_q\) 为映射柱，内含 \(Z=Z\times\{0\}\)，另一端为 \(B\)。\(M\) 形变收缩到 \(B\)，且 \((M,Z)\) 为 CW 对。同伦纤维与相对同伦群的长正合列比较给出
	\[
		\pi_i(M,Z)=0\ (i\ne n+1),\qquad \pi_{n+1}(M,Z)\cong G.
	\]
	具体地，把 \(q\) 作路径纤维化替换后，相对圆盘的基底投影与纤维中的边界互相对应，得到 \(\pi_i(M,Z)\cong\pi_{i-1}(\operatorname{hofib}q)\)（\(i\ge2\)）。这也是第8章纤维化长正合列的空间对形式。

	将 \(Z\) 压成一点，记 \(Q=M/Z\)。相对 Hurewicz 定理用于单连通的 \(M,Z\)，说明 \(H_i(M,Z)=0\)（\(i\le n\)），且 \(H_{n+1}(M,Z)\cong G\)。商的相对胞腔链给出 \(\widetilde H_i(Q)\cong H_i(M,Z)\)。此外 \(Q\) 单连通：把商视作在 \(M\) 上沿 \(Z\) 附着一个锥，van Kampen 将其基本群识别为 \(\pi_1(M)\) 除去 \(\pi_1(Z)\) 的像的正规闭包，两者在这里均为零。依次使用首个非零次数的 Hurewicz 定理，便知 \(Q\) 为 \(n\) 连通，且 \(\pi_{n+1}(Q)\cong G\)。相对与绝对 Hurewicz 映射的自然性还说明，商映射诱导
	\[
		\pi_{n+1}(M,Z)\xrightarrow{\cong}\pi_{n+1}(Q).
	\]

	现在从 \(Q\) 开始附着维数至少为 \(n+3\) 的胞腔，逐次消去更高同伦群，得到 \(K=K(G,n+1)\) 及 \(v:Q\to K\)。复合 \(\bar\kappa:M\to Q\xrightarrow{v}K\) 在 \(Z\) 上恒取基点，并在全部相对同伦群上给出同构
	\[
		\pi_i(M,Z)\xrightarrow{\cong}\pi_i(K,*).
	\]
	将 \(\bar\kappa\) 限制到映射柱的另一端 \(B\)，定义 \(\kappa:B\to K\)。对 \(z\in Z\)，反向走过映射柱线段，得到道路 \(\lambda_z(t)=\bar\kappa([z,1-t])\)，其起点为 \(\kappa(q(z))\)、终点为 \(*\)。因此
	\[
		j(z)=(q(z),\lambda_z)
	\]
	是一张所需的连续映射。

	最后验证弱等价。比较空间对 \((M,Z)\) 的长正合列与 \(E_\kappa\to B\to K\) 的同伦纤维序列，得到
	\[
		\begin{tikzcd}[column sep=small]
			\pi_{i+1}(B)\arrow[r]\arrow[d,equal]&\pi_{i+1}(M,Z)\arrow[r]\arrow[d,"\cong"]&\pi_i(Z)\arrow[r]\arrow[d,"j_*"]&\pi_i(B)\arrow[r]\arrow[d,equal]&\pi_i(M,Z)\arrow[d,"\cong"]\\
			\pi_{i+1}(B)\arrow[r]&\pi_{i+1}(K)\arrow[r]&\pi_i(E_\kappa)\arrow[r]&\pi_i(B)\arrow[r]&\pi_i(K).
		\end{tikzcd}
	\]
	交换性来自 \(\bar\kappa\) 及刚写出的零伦 \(\lambda_z\)。两侧相对群的竖箭头已为同构，故长正合列给出所有 \(j_*\) 为同构；低阶项还说明两空间均单连通。所诱导的纤维比较在 \(\pi_n\) 上是某个 \(G\) 的自同构。改变 \(v\) 在唯一非零同伦群上的识别，即后接实现其逆的 \(K(G,n+1)\) 自同伦等价，就可使这张纤维比较为指定的恒等。这样的自同伦等价可按第11章 Eilenberg–MacLane 模型比较中的逐胞腔构造取得。
\end{proof}

构造中两次处理有不同作用：商 \(M/Z\) 把“在总空间上已经为零”的信息变成一个带基点对象；随后附着高维胞腔，只保留所需的 \(G\)。映射柱线段把这份相对信息重新变成 \(\kappa q\) 的一个具体零伦，所以能恢复总空间到同伦纤维的映射。

\ProseParagraphBreak

为了从这个构造读出上同调类，给 \(K(G,n+1)\) 选普遍类
\[
	\iota_{n+1}\in H^{n+1}(K(G,n+1);G)\cong\operatorname{Hom}(G,G),\qquad
	\iota_{n+1}\longleftrightarrow\operatorname{id}_G.
\]
这个同构来自 Hurewicz 与万有系数定理，低于 \(n+1\) 的约化同调均为零。终点固定的路径空间可缩，所以空间对的连接同态 \(\delta_{\mathrm{in}}\) 给出同构。用
\[
	\sigma_{\mathrm{in}}\iota_{n+1}=\delta_{\mathrm{in}}^{-1}\operatorname{ev}_0^*\iota_{n+1}
\]
规定 \(H^n(\Omega K;G)\) 中的恒等类，即按这个类识别纤维的 \(\pi_n\cong G\)。第11章的相对传递判据直接给出
\[
	\mathrm d_{n+1}(\operatorname{id}_G)=\iota_{n+1}.
\]
这里仅使用加性谱序列；其构造对任意阿贝尔系数成立，不要求 \(G\) 自带环乘法。这个识别也明确记录了道路方向：反向道路映射 \(r:P^{\mathrm{in}}K\to P_*K\) 覆盖 \(K\) 上的恒等，故 \(r|_{\Omega K}^*(\sigma\iota_{n+1})=\sigma_{\mathrm{in}}\iota_{n+1}\)。在同一环路空间中，反向环路在 \(\pi_n(\Omega K)\) 上取负，因此入向与出向的两种悬挂类相差负号。引理中允许调整的纤维识别正是按这里的约定选择。

\begin{definition}[阶段的 \(k\)-不变量]\label{g2:def:postnikov-k-invariant}
	设 \(X\) 为单连通带基点 CW 复形，\(n\ge2\)，记 \(G=\pi_n(X)\)。在阶段纤维化\eqref{g2:eq:postnikov-stage-fibration}中，用 \(\pi_n(F_n)\cong G\) 将纤维类 \(\iota_F\in H^n(F_n;G)\) 识别为 \(\operatorname{id}_G\)。定义
	\begin{equation}\label{g2:eq:postnikov-k-transgression}
		k_{n+1}(X)=\mathrm d_{n+1}\iota_F\in H^{n+1}(P_{n-1}X;G),
	\end{equation}
	称为这一阶段的\term[k bubianliang]{\(k\)-不变量}[\(k\)-invariant]。下标 \(n+1\) 在本节记录上同调次数。
\end{definition}

这个传递在整个目标群中取值。因为 \(F_n\) 为 \((n-1)\) 连通，纤维上同调的第 \(1,\ldots,n-1\) 行全为零，没有更早微分改变 \((0,n)\) 与 \((n+1,0)\) 两项。基底单连通又使 \(G\) 的局部系统为常值。

\begin{theorem}[由 \(k\)-不变量恢复阶段]\label{g2:thm:postnikov-k-reconstruction}
	设 \(X\) 为单连通带基点 CW 复形。对 \(n\ge2\)，令 \(G=\pi_n(X)\)，并用\cref{g2:thm:em-representability}将 \(k_{n+1}(X)\) 表示为映射
	\[
		\kappa_{n+1}:P_{n-1}X\longrightarrow K(G,n+1).
	\]
	则 \(P_nX\) 与 \(\operatorname{hofib}(\kappa_{n+1})\) 具有相同的弱同伦型，比较保持到 \(P_{n-1}X\) 的映射及指定的纤维同伦群识别。在 CW 模型上，这给出同伦等价。特别地，若 \(k_{n+1}(X)=0\)，这一阶段分裂为 \(P_{n-1}X\times K(G,n)\) 的同伦型。
\end{theorem}
\begin{proof}
	先在 \(X[n]\to X[n-1]\) 的 CW 模型上应用\cref{g2:lem:principal-em-fibration}，再经阶段同伦等价搬到 \(P_nX\to P_{n-1}X\)。得到某张 \(\kappa\) 及纤维上诱导恒等的弱等价。比较此纤维化与普遍路径纤维化，传递的自然性给出
	\[
		\mathrm d_{n+1}\iota_F=\kappa^*\iota_{n+1}.
	\]
	左边正是定义的 \(k_{n+1}(X)\)。所以构造出的 \(\kappa\) 表示所需类，可表性定理说明它与任意选定的 \(\kappa_{n+1}\) 同伦。

	还需检查更换代表映射不改变总空间。若 \(H:B\times I\to K\) 是 \(\kappa_0\) 到 \(\kappa_1\) 的同伦，则在式\eqref{g2:eq:k-invariant-path-model}中的道路 \(\lambda\) 前接上反向道路 \(t\mapsto H(b,1-t)\)，得到覆盖 \(B\) 的映射 \(E_{\kappa_0}\to E_{\kappa_1}\)。用原向同伦给出逆；往返道路可保持端点收缩，故两复合与恒等映射纤维同伦。因而同伦类相同的代表给出相同的纤维同伦型。

	当类为零时，取常值代表 \(\kappa=0\)。此时式\eqref{g2:eq:k-invariant-path-model}直接成为 \(B\times\Omega K(G,n+1)\)，其 CW 同伦型即 \(B\times K(G,n)\)。
\end{proof}

可以把这一构造画成拉回方块：
\[
	\begin{tikzcd}[column sep=large,row sep=large]
		P_nX\arrow[r]\arrow[d,"p_n"']&P^{\mathrm{in}}K(G,n+1)\arrow[d,"\operatorname{ev}_0"]\\
		P_{n-1}X\arrow[r,"\kappa_{n+1}"']&K(G,n+1).
	\end{tikzcd}
\]
在本图中用同伦等价的模型替换 \(P_nX\)，上边的路径空间可缩，纤维均为 \(K(G,n)\)。底边的一张映射便决定了这一层怎样附着。

\ProseParagraphBreak

\(k_{n+1}\) 与阶段纤维化的首截面障碍相联系，但须遵守两种边界的方向。若障碍 \(o_{n+1}\) 按胞腔正向边界上的截面映射取值，则本书的约定给出
\[
	o_{n+1}=-k_{n+1}.
\]
可在单个胞腔上看见这个符号。把纤维化拉回 \(D^{n+1}\)，在圆盘上选一个局部截面，其边界记为 \(l\)；既定骨架截面的边界记为 \(s\)。两者的差 \(s-l\) 经纤维方向的变形代表障碍 \(o_{n+1}\)。因此与 \(\iota_F\) 配对，若 \(A\) 是骨架上延拓纤维恒等类的上圈，就有 \(A(s)-A(l)=o_{n+1}\)。延拓 \(A\) 为下一层的上链后，传递上链在该胞腔上取值 \(A(l)\)；而 \(A(s)\) 是把 \(A\) 沿骨架截面拉回后所得上链的胞腔余边界。模去余边界，便得到 \(\tau\iota_F=-o_{n+1}\)。这与第11章“延拓后取余边界”的传递约定相容。

\ProseParagraphBreak

对给定 \(f:Y\to P_{n-1}X\)，式\eqref{g2:eq:k-invariant-path-model}直接说明：把 \(f\) 提升到 \(P_nX\)，就是给 \(\kappa_{n+1}f\) 选一个零伦；存在条件为
\begin{equation}\label{g2:eq:postnikov-lift-obstruction}
	f^*k_{n+1}(X)=0\in H^{n+1}(Y;G)
\end{equation}
其中 \(Y\) 为 CW 复形。这一次条件是完整条件，因为纤维 \(K(G,n)\) 没有其他非零同伦群。

\ProseParagraphBreak

一个 \(k\)-不变量的具体表达还依赖于对阶段与系数群的识别。若 \(h:B\to B'\) 是基底比较，\(\alpha:G\to G'\) 是系数同构，则对应的相容关系为 \(\alpha_*k=h^*k'\)。系数变化改变纤维的命名，基底比较则改变类所处的空间；恢复阶段时两者都要同时记录。对非单连通空间，逐次消去高阶同伦群仍产生阶段，但 \(\pi_1\) 对 \(\pi_n\) 的作用给出局部系统 \(\Pi_n\)，相应障碍应取在 \(H^{n+1}(P_{n-1}X;\Pi_n)\) 中。作用非平凡时，不能把它替换为上述到普通 \(K(G,n+1)\) 的常系数映射。

\subsection{同伦群相同的不同三型}
\label{g2:subsec:1203:04}

现在给出同伦群完全相同、同伦型却不同的一族空间。取
\[
	B=K(\mathbb Z,2)=\mathbb CP^\infty,\qquad H^*(B;\mathbb Z)=\mathbb Z[u],\quad |u|=2.
\]
对整数 \(m\)，可表性定理给出代表 \(m u^2\in H^4(B;\mathbb Z)\) 的映射
\[
	\kappa_m:B\longrightarrow K(\mathbb Z,4).
\]
令 \(E_m\) 为其同伦纤维的 CW 逼近。等价的路径模型给出纤维化 \(K(\mathbb Z,3)\to E_m\to K(\mathbb Z,2)\)，以下计算在该模型上进行，CW 逼近不改变同伦群与奇异同调。

\begin{example}[相同同伦群的不同三型]\label{g2:exm:postnikov-two-stage-family}
	上述 \(E_m\) 均单连通，并且
	\[
		\pi_i(E_m)\cong
		\begin{cases}
			\mathbb Z,&i=2,3,\\
			0,&i\ne2,3\quad(i\ge1).
		\end{cases}
	\]
	然而
	\[
		H^3(E_m;\mathbb Z)\cong\ker(\mathbb Z\xrightarrow{\times m}\mathbb Z),\qquad
		H^4(E_m;\mathbb Z)\cong\operatorname{coker}(\mathbb Z\xrightarrow{\times m}\mathbb Z).
	\]
	因此不同正整数 \(m\) 给出两两不同的同伦型；\(m=0\) 与 \(m=1\) 也不同。
\end{example}
\begin{proof}
	对 \(E_m\to K(\mathbb Z,2)\to K(\mathbb Z,4)\) 用同伦长正合列。二次群由基底给出；三次群由连接同态 \(\pi_4(K(\mathbb Z,4))\to\pi_3(E_m)\) 给出；其余正阶项为零。这个计算与 \(m\) 无关，因为 \(\kappa_m\) 在全部同伦群上诱导零映射：两端唯一非零的次数不同。

	上同调能看见 \(m\)。纤维 \(K(\mathbb Z,3)\) 可由 \(S^3\) 附着维数至少为五的胞腔构造，故其 \(H_3=\mathbb Z\)、\(H_4=0\)；万有系数定理给出 \(H^3=\mathbb Z\)、\(H^4=0\)，一次、二次也为零。按入向路径纤维化的悬挂类选择三次生成元 \(v\)。在总次数不超过四的范围，Serre 页只有基底的 \(1,u,u^2\) 和 \((0,3)\) 处的 \(v\)。唯一可能的微分为
	\[
		\mathrm d_4:\mathbb Zv\longrightarrow\mathbb Zu^2,\qquad \mathrm d_4v=m u^2,
	\]
	其中等式来自路径模型的自然性。第三次因此留下这张映射的核，第四次留下其余核；每个次数只有一层，故无加性扩张。

	当 \(m=0\) 时，\(E_0\) 为 \(K(\mathbb Z,2)\times K(\mathbb Z,3)\) 的同伦型，三次、四次上同调均为 \(\mathbb Z\)。当 \(m=1\) 时，两群均为零；当 \(m>1\) 时，第四次群为 \(\mathbb Z/m\)。这些群区分了所列空间。
\end{proof}

这个例子也解释 Whitehead 定理中的映射条件。每个 \(E_m\) 都有同样的抽象同伦群，但若一张映射在全部同伦群上诱导同构，它就必须是 CW 模型之间的同伦等价，因而还必须保持刚算出的上同调。上述已由上同调区分的空间之间不存在这样的实现映射；障碍正包含在 \(k_4=m u^2\) 中。

\ProseParagraphBreak

球面 \(S^2\) 给出 \(m=1\) 的一个具体来源。已有 \(\pi_2(S^2)=\pi_3(S^2)=\mathbb Z\)，所以它的第三阶段位于
\[
	K(\mathbb Z,3)\longrightarrow P_3S^2\longrightarrow K(\mathbb Z,2).
\]
\(S^2\to S^2[3]\) 只添加维数至少为五的胞腔。因而 \(H_2=\mathbb Z\)、\(H_3=H_4=0\)；即使五胞腔能对四次同调增加边界关系，原来的四次群已经为零。万有系数定理遂给出 \(H^3(P_3S^2)=H^4(P_3S^2)=0\)。写 \(k_4(S^2)=m u^2\)，上一计算迫使乘 \(m\) 无核且无余核，故 \(m=\pm1\)。改变纤维生成元的符号后可取 \(m=1\)，所以 \(P_3S^2\) 是 \(\operatorname{hofib}(u^2:\mathbb CP^\infty\to K(\mathbb Z,4))\) 的 CW 同伦型。Hopf 类提供三次同伦群，\(u^2\) 则记录这层怎样接在二次阶段上。

\ProseParagraphBreak

对阶段构造、主纤维化与 \(k\)-不变量的进一步阅读，可参看\cite[第4.3节，Postnikov Towers]{Hatcher2002}及\cite[第22章第4节，Postnikov systems]{May1999}。本节采用上同调次数作为 \(k\) 的下标；比较不同文献时，先核对该类所在的 \(H^{n+1}(P_{n-1}X;\pi_nX)\)，再比较下标。
% 纲要条目排为编号小节标题，并进入目录与书签；说明段仍保留于纲要注释中。

% !TeX root = ../../Book2.tex
% 纲要标题：### 12.4 几何中的截面障碍
% 纲要位置：几何学计划纲要.md，第 1493 行。
% 正文状态：初稿；尚未正式审核。

% 纲要内容开始
% 
% * 球丛的首障碍与 Euler 类
% * 向量场及结构群约化
% * 平面丛与曲面上的向量场
% 
% ---
% 
% 纲要内容结束

\section{几何中的截面障碍}
\label{g2:sec:1204}

给每个点连续选一个单位向量，是一个截面问题。沿一个胞腔的边界走一圈，已经选好的向量可能在纤维球面上绕出非零同伦类，这就阻止它填满胞腔。上一章则从纤维定向类的上同调延拓定义了 Euler 类。本节把这两种延拓放到同一个胞腔上比较，再用向量场、平面丛和曲面上的计算说明：障碍类的次数、系数以及消失后的结论都由纤维决定。

\subsection{球丛的首障碍与 Euler 类}
\label{g2:subsec:1204:01}
设 \(p:E\rightarrow B\) 是定向 \(S^m\) 丛，\(m\ge1\)，其中 \(B\) 为道路连通 CW 复形。这里的定向仍指沿道路运输保持纤维的正向顶次类。\cref{g2:prop:sphere-connectivity}给出 \(\pi_j(S^m)=0\) 对 \(1\le j<m\) 成立，再加上纤维道路连通，故可以从零骨架开始逐胞腔选截面，一直得到 \(s:B^m\rightarrow E_m\)，其中 \(E_m=p^{-1}(B^m)\)。第一次可能失败发生在 \((m+1)\) 胞腔：将丛拉回其特征圆盘并平凡化，边界截面成为 \(a:S^m\rightarrow S^m\)。\cref{g2:lem:sphere-degree-classification}说明它零伦当且仅当 \(\deg a=0\)，所以延拓问题成为一个整数是否为零的问题。

\ProseParagraphBreak

源球面取胞腔圆盘的正向边界，目标球面取纤维定向。\cref{g2:thm:primary-obstruction}于是给出与初始截面选择无关的首障碍
\[
	o_{m+1}(p)\in H^{m+1}(B;\mathbb Z).
\]
系数为普通整数，是因为保持纤维定向的运输在 \(\pi_m(S^m)\cong\mathbb Z\) 上作用恒等。\cref{g2:def:spherical-euler-class}所定义的 \(e(p)=\mathrm d_{m+1}u\) 也在这个群里；其中 \(u\) 是纤维的正向上同调生成元。比较二者时需要保留连接同态的符号，下面的局部计算同时给出识别及其符号。

\begin{theorem}[球丛的首障碍]\label{g2:thm:sphere-primary-euler}
	设 \(p:E\rightarrow B\) 为定向 \(S^m\) 丛，\(m\ge1\)，\(B\) 为道路连通 CW 复形。障碍上链按正向胞腔边界上的截面映射度取值，Euler 类按式\eqref{g2:eq:spherical-euler-transgression}定义，则
	\begin{equation}\label{g2:eq:primary-euler-sign}
		o_{m+1}(p)=-e(p).
	\end{equation}
	若 \(\dim B\le m+1\)，则 \(p\) 有连续截面当且仅当 \(e(p)=0\)。对定向圆周丛 \(m=1\)，这一等价对任意维数的 CW 底空间均成立。
\end{theorem}
\begin{proof}
	先比较同一个胞腔上的两种延拓。固定 \(B^m\) 上的截面 \(s\)。由\cref{g2:prop:serre-transgression}，纤维生成元 \(u\) 可延拓为 \(v\in H^m(E_m;\mathbb Z)\)：底空间只有 \(m\) 维，球丛谱序列中唯一可能离开 \((0,m)\) 的微分已落在其维数以外。以
	\[
		v-p^*s^*v
	\]
	替换 \(v\)，不改变它在纤维上的限制，并使 \(s^*v=0\)。减去这一底空间类，是为了把截面所带的贡献从延拓类中消去。

	取一个正向 \((m+1)\) 胞腔的特征映射 \(\chi:D^{m+1}\rightarrow B\)，在拉回丛 \(D^{m+1}\times S^m\) 中使用保持纤维定向的平凡化。令 \(x\in H^m(\partial D^{m+1};\mathbb Z)\) 为正向边界生成元。边界上的乘积有
	\[
		H^m(\partial D^{m+1}\times S^m;\mathbb Z)=\mathbb Z\operatorname{pr}_1^*x\oplus\mathbb Z\operatorname{pr}_2^*u.
	\]
	\(v\) 在纤维上的值为 \(u\)，所以它在这个边界乘积上的限制必为
	\[
		\operatorname{pr}_2^*u+\lambda\operatorname{pr}_1^*x,\qquad \lambda\in\mathbb Z.
	\]
	边界截面是图映射 \(\Gamma_a(z)=(z,a(z))\)。把上式沿 \(\Gamma_a\) 拉回，条件 \(s^*v=0\) 给出
	\[
		0=a^*u+\lambda x=(\deg a+\lambda)x,
	\]
	从而 \(\lambda=-\deg a\)。

	现在使用空间对 \((D^{m+1}\times S^m,\partial D^{m+1}\times S^m)\) 的连接同态。纤维类 \(\operatorname{pr}_2^*u\) 已定义在整个圆盘乘积上，故连接同态将它送到零；\(\operatorname{pr}_1^*x\) 则被送到圆盘的正向相对生成元乘纤维的零次单位元。后一个正号可以直接由上链检查：若 \(\widetilde x\) 延拓边界代表元，则
	\[
		\langle\mathrm d\widetilde x,[D^{m+1}]\rangle=\langle x,\partial[D^{m+1}]\rangle=1.
	\]
	因此 \(v\) 的上同调延拓障碍在该胞腔上的值是 \(\lambda=-\deg a\)。\cref{g2:prop:serre-transgression}把这些值组成的上圈识别为 \(\mathrm d_{m+1}u\)，而截面首障碍的值为 \(\deg a\)。逐胞腔比较后得到式\eqref{g2:eq:primary-euler-sign}。

	若 \(\dim B\le m+1\)，首障碍为零时，\cref{g2:thm:primary-obstruction}允许先修改 \(B^m\) 上的截面，再将它延拓到全部 \((m+1)\) 胞腔；这已经覆盖整个 \(B\)。反之，全局截面使全部胞腔障碍为零。结合比较式便得到所述等价。若 \(m=1\)，首障碍为零后先延拓到 \(B^2\)，以后每一步的系数群 \(\pi_q(S^1)\)（\(q\ge2\)）都为零，因而可以继续延拓。各骨架上的截面相容，由 CW 弱拓扑粘合为全局连续截面。
\end{proof}

式\eqref{g2:eq:primary-euler-sign}中的负号来自已经固定的两种约定：截面障碍取边界映射的正向度，上一章 Euler 类取“延拓后取余边界”的传递。若采用首障碍来规范 Euler 类，其记号对应这里的 \(-e(p)\)。判定截面存在时只使用这个类是否为零，因而两种规范给出相同判定；计算它在定向基本类上的具体整数时则必须保持所选符号。

\ProseParagraphBreak

对 \(m\ge2\)，\(e(p)=0\) 保证可以选到 \(B^{m+1}\) 上的截面。继续向上延拓仍要考察 \(\pi_{m+1}(S^m),\pi_{m+2}(S^m),\ldots\) 中的障碍，下一步落在 \(H^{m+2}(B;\Pi_{m+1}(p))\) 中，并可能依赖已选截面。球面低阶同伦群的消失给出上述维数范围；圆周的全部高阶同伦群消失，则使其首障碍成为任意维数上的完整判据。若球丛不定向，首障碍仍可构造，但系数须换为由纤维定向反转给出作用的整数局部系统；上述普通整数 Euler 类的陈述以定向为条件。

\subsection{向量场及结构群约化}
\label{g2:subsec:1204:02}
设 \(V\rightarrow B\) 为实秩 \(r\) 向量丛，并给定一个连续正定内积。记 \(S(V)\) 为各纤维单位球面组成的球丛，\(\varepsilon_B^k=B\times\mathbb R^k\) 为平凡秩 \(k\) 丛。无零截面 \(v\) 可逐点归一化为 \(v/\|v\|\)，所以寻找无零截面就是寻找 \(S(V)\) 的截面。对光滑流形上的光滑向量丛，第一卷第12章第3节的内积存在定理保证了这样的光滑内积；其假设为流形有限维、Hausdorff 且第二可数，证明是用局部有限的光滑单位分解将局部 Euclid 内积相加。

\ProseParagraphBreak

还可以从标架看这个问题。局部选取正交标架后，过渡矩阵属于 \(\operatorname O(r)\)。若第一根标架向量能够在不同图上保持不变，过渡矩阵便有 \(\operatorname{diag}(1,A)\) 的形式。这提示我们，整体向量场对应于过渡矩阵所能取值的群变小。

\begin{definition}[正交结构群约化]\label{g2:def:orthogonal-reduction}
	设 \(V\rightarrow B\) 是带连续正定内积的实秩 \(r\) 向量丛，\(0\le k\le r\)。若存在局部正交标架，使所有过渡矩阵都落在子群
	\[
		\left\{\begin{pmatrix}I_k&0\\0&A\end{pmatrix}:A\in\operatorname O(r-k)\right\}\subseteq\operatorname O(r),
	\]
	则称 \(V\) 的正交\term[jiegouqun-yuehua]{结构群约化}[reduction of structure group]到 \(\operatorname O(r-k)\)。这里约定 \(\operatorname O(0)\) 为平凡群。对定向丛，使用正向正交标架，并将两个正交群换成相应的特殊正交群，就得到到 \(\operatorname{SO}(r-k)\) 的约化。
\end{definition}

\begin{proposition}[独立截面与约化]\label{g2:prop:sections-orthogonal-reduction}
	设 \(V\rightarrow B\) 为带连续正定内积的实秩 \(r\) 向量丛，\(1\le k\le r\)。以下条件等价：
	\begin{equivalences}
		\item \(V\) 有 \(k\) 个逐点线性无关的连续截面；
		\item \(V\) 有 \(k\) 个逐点正交归一的连续截面；
		\item 存在秩 \(r-k\) 向量丛 \(W\)，使 \(V\cong\varepsilon_B^k\oplus W\)；
		\item 正交结构群可约化到 \(\operatorname O(r-k)\)。
	\end{equivalences}
	对光滑向量丛及光滑内积，同一等价在全部截面、丛同构和局部标架都要求光滑时成立。若 \(V\) 已定向，则上述条件也等价于到 \(\operatorname{SO}(r-k)\) 的约化。
\end{proposition}
\begin{proof}
	对逐点独立的截面 \(v_1,\ldots,v_k\) 作 Gram--Schmidt 正交化：令
	\[
		w_j=v_j-\sum_{i<j}\langle v_j,e_i\rangle e_i,\qquad e_j=\frac{w_j}{\|w_j\|}.
	\]
	独立性保证每个 \(w_j\) 都处处非零，故这些公式保持连续性；在光滑情形也保持光滑性。由此得到 I 推出 II，反向蕴含直接成立。

	给定 II 中的正交截面，令 \(W_b\) 为它们在 \(V_b\) 中的共同正交补。映射
	\[
		P_b(v)=v-\sum_{i=1}^k\langle v,e_i(b)\rangle e_i(b)
	\]
	是到 \(W_b\) 的投影。在某点选取 \(r-k\) 个局部截面，其投影在该点构成 \(W_b\) 的基；相应 Gram 行列式在附近仍非零，因此这些投影给出 \(W\) 的局部标架。这证明 \(W\) 确为连续子丛，光滑时也为光滑子丛。映射
	\[
		\varepsilon_B^k\oplus W\longrightarrow V,\qquad (b,t_1,\ldots,t_k,w)\longmapsto\sum_i t_i e_i(b)+w
	\]
	是丛同构，其逆由内积系数与 \(P_b\) 给出，故 II 推出 III。一个直和分解中平凡因子的标准截面又给出 I。

	在 \(W\) 上局部正交化标架，将它接到 \(e_1,\ldots,e_k\) 后面，便得到 \(V\) 的局部正交标架。不同图上的前 \(k\) 根相同，过渡矩阵因而具有所述块对角形式，得到 IV。反过来，这种过渡矩阵固定前 \(k\) 根标准向量，所以各图中的这些局部截面粘成全局正交截面。定向情形下，当 \(k<r\) 时，由 \(V\) 的定向及有序截面组赋予 \(W\) 相容定向，再使用 \(W\) 的正向正交标架。当 \(k=r\) 时，可在每个连通分支上按需反转最后一根向量，使全标架保持给定定向。于是同样得到特殊正交群的结论。
\end{proof}

这个证明使用的正交补构造也见第一卷第12章第3节；这里把投影写出，是为了说明它同时适用于连续丛，以及约化中的标架如何实际产生。第一卷第12章第1节的“整体标架与平凡性”正是上述 \(k=r\) 的情形。

\ProseParagraphBreak

把一个纤维中全部有序正交 \(k\) 标架放在一起，得到\term[stiefel-cong]{Stiefel 丛}[Stiefel bundle] \(V_k(V)\rightarrow B\)，其纤维是
\[
	V_k(\mathbb R^r)=\{(e_1,\ldots,e_k):\langle e_i,e_j\rangle=\delta_{ij}\}.
\]
局部正交标架将它平凡化为 \(U\times V_k(\mathbb R^r)\)，过渡矩阵同时作用于每根向量，所以此构造与局部选择相容。它的截面就是 \(k\) 个正交单位截面。\(k=1\) 恢复单位球丛；\(k>1\) 时，延拓所用的纤维同伦群变成 Stiefel 空间的同伦群。这将“找多少个独立向量场”统一成了截面问题。\symidx{\(V_k(\mathbb R^r)\)：有序正交 \(k\) 标架的 Stiefel 空间}

\ProseParagraphBreak

在 \(r\) 维光滑流形 \(M\) 上取 \(V=\mathrm TM\)，处处非零向量场等价于单位切丛的截面，也等价于正交结构群从 \(\operatorname O(r)\) 约化到 \(\operatorname O(r-1)\)。这里连续存在性与光滑存在性一致。事实上，给一个连续无零截面 \(v\)，在每点附近取局部光滑截面 \(w_i\)，使 \(\|w_i-v\|<\|v\|/2\)；可以在局部标架中固定该点的系数，再由连续性缩小邻域来做到。用第一卷第12章第3节的光滑单位分解 \(\{\varphi_i\}\) 形成 \(w=\sum_i\varphi_iw_i\)，便有
\[
	\|w-v\|\le\sum_i\varphi_i\|w_i-v\|<\frac{\|v\|}{2},
\]
故 \(w\) 是光滑无零截面。在维数不超过 \(r\) 的 CW 底空间上，\(r\ge2\) 的定向秩 \(r\) 向量丛只有首个截面障碍落在可见维数内，所以 \(e(V)=0\) 正好判定无零截面的存在。

\ProseParagraphBreak

秩一时单位纤维是离散的 \(S^0\)，延拓首先要求两端所选分支能够连起来。例如 Möbius 线丛
\[
	L=([0,1]\times\mathbb R)/\bigl((1,t)\sim(0,-t)\bigr)\longrightarrow S^1
\]
的截面在区间上由连续函数 \(a\) 表示，端点条件为 \(a(1)=-a(0)\)。若 \(a(0)=0\)，已经有零点；若不为零，介值定理仍迫使 \(a\) 在区间内取零。单位纤维的两个点绕基底一圈后交换，正是无法选取连续分支的原因。这个障碍发生在一维，属于离散纤维的分支问题。

\subsection{平面丛与曲面上的向量场}
\label{g2:subsec:1204:03}
二维底空间配圆周纤维时，整个延拓过程只需一个整数上圈：先在顶点和边上选截面，再读取各面边界的绕数。下面先让过渡函数直接给出这个整数，再在切丛上把它拆成两个局部贡献。

\begin{example}[球面上的旋转粘合]\label{g2:exm:clutched-plane-obstruction}
	给 \(S^2\) 一个定向，把它写成两张闭圆盘 \(D_+\cup D_-\)。在公共边界上用 \(t\in\mathbb R/2\pi\mathbb Z\) 表示 \(D_-\) 的正向边界参数。对整数 \(k\)，将两张平凡定向实平面丛按
	\[
		(t,z)_+\sim(t,e^{\mathrm i kt}z)_-,\qquad z\in\mathbb C\cong\mathbb R^2,
	\]
	粘合，得到定向平面丛 \(V_k\rightarrow S^2\)；Euclid 内积因旋转粘合而给出全局内积。

	在 \(D_+\) 上选常值单位截面 \(1\)，它在 \(D_-\) 的边界坐标中变成 \(e^{\mathrm i kt}\)，绕数为 \(k\)。把 \(D_+\) 作为已经固定截面的子复形，填入 \(D_-\) 的相对首障碍就是这个整数。由于 \(D_+\) 可缩，\(H^2(S^2,D_+;\mathbb Z)\rightarrow H^2(S^2;\mathbb Z)\) 为同构；正向 \(D_-\) 对应满足 \(\langle\omega,[S^2]\rangle=1\) 的生成元 \(\omega\)。于是
	\[
		o_2(S(V_k))=k\omega,\qquad e(V_k)=-k\omega.
	\]
	\cref{g2:thm:sphere-primary-euler}说明 \(V_k\) 有无零截面当且仅当 \(k=0\)。这里即使允许修改 \(D_+\) 上的截面，结论也相同：其沿边界的附加绕数为零，因为它已经延拓到 \(D_+\) 内部。\(k=0\) 时粘合为恒等，两个标准基向量直接粘成整体标架。
\end{example}

对闭定向连通曲面 \(\Sigma_g\) 上的定向平面丛 \(V\)，同一个方法把存在性归结为一个整数。\cref{g2:prop:surface-integer-homology,g2:thm:uct-cohomology}给出 \(H^2(\Sigma_g;\mathbb Z)\cong\mathbb Z\)；取正向生成元 \(\omega\)，写 \(e(V)=q\omega\)。于是 \(V\) 有无零截面当且仅当 \(q=0\)。一旦选定一个单位截面 \(s_0\)，平面定向与内积给出逐纤维的正向四分之一转 \(J\)，\((s_0,Js_0)\) 就是全局正向正交标架。因此每个单位截面都唯一写为
\[
	s=f_1s_0+f_2Js_0,\qquad f=(f_1,f_2):\Sigma_g\longrightarrow S^1.
\]
截面的同伦对应这些映射的同伦。\(S^1\) 是 \(K(\mathbb Z,1)\)，故\cref{g2:thm:em-representability}给出
\[
	\{\text{单位截面的同伦类}\}\cong H^1(\Sigma_g;\mathbb Z)\cong\mathbb Z^{2g}.
\]
这里自由同伦与保基点同伦给出相同的分类：选基点 \(b_0\)，将 \(f\) 换成 \(f/f(b_0)\) 就使基点映到 \(1\)，而任意同伦 \(H\) 也可规范为 \(H(b,t)/H(b_0,t)\)。这个坐标依赖于 \(s_0\)：更换参考截面会平移所有坐标。各整数可由曲面标准 \(2g\) 条生成回路上 \(f\) 的绕数读出；曲面二胞腔的附着字是交换子的乘积，其总绕数为零，所以这些整数可以任意指定。

\begin{example}[球面与环面的切向量场]\label{g2:exm:tangent-sphere-torus-obstruction}
	给 \(S^2\subset\mathbb R^3\) 外向定向，令 \(e_3=(0,0,1)\)。将常向量 \(e_3\) 正交投影到切平面，得到光滑切向量场
	\[
		v(x)=e_3-x_3x.
	\]
	因为 \(\langle v(x),x\rangle=x_3-x_3=0\)，它确实切于球面；其零点恰为南北极 \(S=-e_3,N=e_3\)。删去两极附近的互不相交小圆盘，余下的圆环记为 \(A\)。在 \(A\) 上归一化 \(v\)，就有单位切丛的截面。要填回两个圆盘，分别计算这两个边界截面的绕数。

	在北极附近用 \((x_1,x_2)\) 作坐标，\(x_3=\sqrt{1-x_1^2-x_2^2}\)。向量场在坐标切标架中的两个系数是
	\[
		(-x_3x_1,-x_3x_2).
	\]
	沿小圆周把正函数 \(x_3\) 连续变到 \(1\)，该非零系数向量同伦到 \((-x_1,-x_2)\)，其归一化是圆周上的旋转半圈，度为 \(+1\)。南极附近 \(x_3=-\sqrt{1-x_1^2-x_2^2}\)，同样的系数向量同伦到 \((x_1,x_2)\)，故度也为 \(+1\)。南极这张坐标与外向定向相反，但改用正向坐标会同时反转圆盘参数及坐标切标架的定向，两次符号变化相消，所以局部贡献仍为 \(+1\)。

	取以 \(A\) 为子复形、两张正向圆盘为剩余二胞腔的相对 CW 结构。这两个障碍组成 \(H^2(S^2,A;\mathbb Z)\cong\mathbb Z^2\) 中的 \((1,1)\)。忘掉固定在 \(A\) 上的截面，得到绝对首障碍；相对类到绝对类的映射将两个值相加，因为正向基本圈在两张圆盘上的系数都是 \(1\)。因此
	\[
		o_2(S(\mathrm TS^2))=2\omega,\qquad e(\mathrm TS^2)=-2\omega
	\]
	沿用式\eqref{g2:eq:primary-euler-sign}的规范。这里零点附近正向小圆周上的归一化向量场之度，就是该零点的\term[lingdian-jubu-zhibiao]{局部指标}[local index]；两极的指标总和为 \(2\)，对应首障碍 \(o_2\) 的取值。这证明球面不存在处处非零连续切向量场，也解释了为什么改变原向量场的局部形状不能把所有零点一起消去：所得绝对障碍类保持不变。

	环面 \(T^2=\mathbb R^2/(2\pi\mathbb Z)^2\) 则有两个整体向量场 \(\partial/\partial\theta_1\)、\(\partial/\partial\theta_2\)。它们来自 \(\mathbb R^2\) 的常向量场，在格点平移下不变，所以可下降，且处处独立。故 \(\mathrm TT^2\cong\varepsilon_{T^2}^2\)，首障碍及 Euler 类均为零；选择这个标架后，单位向量场的同伦类由两条基本回路上的绕数 \((a,b)\in\mathbb Z^2\) 完全描述。
\end{example}

这些计算分出了三个不同层次：Möbius 线丛在选取单位纤维的分支时失败；定向平面丛的分支问题自动解决，首障碍是面边界的整数绕数；首障碍消失之后，已有截面的不同同伦类还由低一次的上同调记录。截面是否存在与存在之后有多少种同伦选择，因而分别对应本章的障碍类与差类。
% 纲要条目排为编号小节标题，并进入目录与书签；说明段仍保留于纲要注释中。


\begin{chaptersummary}
	单个胞腔上的延拓取决于边界球面映射是否零伦；沿基底运输这些同伦类，便得到障碍所用的局部系数。逐胞腔障碍组成余圈，改变上一层选择的差异上链满足 \(\delta D=c(t)-c(s)\)。因此上链为零判定当前截面能够延拓，上同调类为零则允许先修改上一骨架。纤维在首个非零同伦群以下连通时，首障碍与这些初始选择无关；延拓存在后，在维数或纤维同伦群的相应范围内，差异类使同伦类组成上同调群的主齐性集。以 \(K(G,n)\) 为目标便得到 \(H^n(-;G)\) 的可表性。单连通空间的 Postnikov 塔逐次保留低阶同伦群，阶段间的纤维为 \(K(\pi_nX,n)\)，\(k_{n+1}\) 则决定这一纤维如何接上去。几何中，单位向量是球丛截面，多个独立向量对应正交结构群约化；首截面障碍与传递定义的 Euler 类按本书约定相差负号，消失性给出低维底空间上的完整存在判据。
\end{chaptersummary}

\BookExercises

\begin{exercise}\label{g2:ex:torus-extension-commutator}
	将环面 \(T^2\) 表示为在 \(S^1\vee S^1\) 上沿交换子附着一个二胞腔。给定整数 \(a,b\)，在两个圆周上分别取绕数为 \(a,b\) 的映射到 \(S^1\)，证明它们可以延拓到整个环面，并写出一个显式延拓。若把目标换为 \(S^1\vee S^1\)，让两个圆周分别映到目标的两个标准生成圆周，则能否延拓？比较两个问题的区别。
\end{exercise}
\begin{solution}
	设环面的两条一胞腔为 \(u,v\)，二胞腔沿 \(uvu^{-1}v^{-1}\) 附着。目标 \(S^1\) 的基本群是交换群 \(\mathbb Z\)，所以边界复合的同伦类为 \(a+b-a-b=0\)，单胞腔延拓判据给出延拓。将环面写为 \(S^1\times S^1\)，映射
	\[
		f(z,w)=z^a w^b
	\]
	就是一个显式选择；在 \(S^1\times\{1\}\) 与 \(\{1\}\times S^1\) 上分别恢复给定幂映射。如果原给定映射只与幂映射有相同绕数，先在一骨架上取相应同伦，再用 CW 同伦延拓性质把末端边界换回原给定映射。

	目标 \(S^1\vee S^1\) 的基本群是两个生成元 \(x,y\) 的自由群。二胞腔边界的像为 \(xyx^{-1}y^{-1}\)，这是非空约化字，故不是单位元，不能延拓。两个问题中的边界在一次同调中都为零，因为同调只记录交换化；能否填圆盘却由实际基本群中的类决定。非交换的一次障碍因而不能只用整数绕数或普通一次同调替代。
\end{solution}

\begin{exercise}\label{g2:ex:reflection-circle-sections}
	令
	\[
		E=(S^1\times[0,1])/((z,1)\sim(\bar z,0))\longrightarrow S^1
	\]
	为以复共轭粘合的圆周丛。证明它有截面，并直接计算截面在纤维内同伦下的类数。再算 \(H^1(S^1;\Pi_1(p))\)，比较两个答案。
\end{exercise}
\begin{solution}
	常值纤维坐标 \(1\) 与 \(-1\) 都满足粘合条件，因此给出截面。任一截面可写成 \(t\mapsto[u(t),t]\)，其中 \(u:[0,1]\to S^1\) 满足 \(u(0)=\overline{u(1)}\)。区间可缩，故可取实值辐角提升 \(u(t)=e^{\mathrm i\theta(t)}\)。端点条件等价于
	\[
		\theta(0)+\theta(1)=2\pi k\qquad(k\in\mathbb Z).
	\]
	改变辐角提升会使 \(\theta\) 整体增加 \(2\pi\ell\)，从而把 \(k\) 改为 \(k+2\ell\)。因此 \(k\) 的奇偶性与提升选择无关。

	若有截面同伦，则相应的 \(u(t,s)\) 在方形上可取连续实值辐角提升。\(\bigl(\theta(0,s)+\theta(1,s)\bigr)/(2\pi)\) 是随 \(s\) 连续的整数，故保持不变，说明奇偶性为同伦不变量。反过来，先用线性插值把 \(\theta(t)\) 变为连接原两端值的仿射函数，再让初端值从 \(\theta(0)\) 移到 \(\pi k\)，末端值始终取 \(2\pi k\) 减初端值。这些同伦都保留端点条件，最后得到常值函数 \(\pi k\)。所以每张截面同伦于常值纤维坐标 \((-1)^k\)，且这两类不同，恰有两类。

	复共轭在 \(\pi_1(S^1)=\mathbb Z\) 上作用为乘 \(-1\)，故 \(\Pi_1(p)\) 是反号局部系统。圆周的局部系数胞腔上链为 \(\mathbb Z\xrightarrow{-2}\mathbb Z\)，因此 \(H^1(S^1;\Pi_1(p))=\mathbb Z/2\)。这与两类截面的差异一致；若误用常值系数，就会得到 \(\mathbb Z\)，错误地预测无穷多个差异类。
\end{solution}

\begin{exercise}\label{g2:ex:homology-invisible-extension-obstruction}
	考虑 Hopf 映射 \(h:S^3\to S^2\)。证明它在全部正次整数同调上诱导零映射，却不能延拓到 \(D^4\)。说明这与第9章 Hurewicz 定理的范围如何相容。
\end{exercise}
\begin{solution}
	\(S^3\) 的正次同调只在三次为 \(\mathbb Z\)，但 \(H_3(S^2)=0\)，因此所有正次诱导同调映射都是零。由\cref{g2:thm:hopf-infinite-order}，Hopf 映射在 \(\pi_3(S^2)\) 中具有无穷阶，尤其非零。单胞腔延拓判据遂说明它不能填成 \(D^4\to S^2\)。

	\(S^2\) 是一连通空间，Hurewicz 定理保证首个非零次数的 \(\pi_2(S^2)\to H_2(S^2)\) 为同构，并不把三次同伦群识别为三次同调。这里的障碍正位于下一次的 \(\pi_3(S^2)\)，所以同调看不见它并不违反该定理。延拓问题必须保留边界所在的具体同伦次数。
\end{solution}

\begin{exercise}\label{g2:ex:obstruction-cochain-versus-class}
	设 \(m\ge2\)，\(X=S^1\cup_m e^2\) 为沿度 \(m\) 映射附着二胞腔的空间。在一骨架上给定 \(s_k:S^1\to S^1\)、\(s_k(z)=z^k\)。计算其延拓障碍上链及在 \(H^2(X;\mathbb Z)\) 中的类，并判断保留整个一骨架或仅保留顶点时分别能否延拓。最后把固定一骨架的问题写成相对障碍类。
\end{exercise}
\begin{solution}
	取胞腔定向使二维边界为乘 \(m\)。附着圆周经 \(s_k\) 后的度为 \(mk\)，所以障碍上链是 \(c^2(s_k)=mk\in C^2(X;\mathbb Z)=\mathbb Z\)。绝对胞腔上链中 \(\delta^1:\mathbb Z\to\mathbb Z\) 为乘 \(m\)，故 \(H^2(X;\mathbb Z)=\mathbb Z/m\)，而 \([c^2(s_k)]=0\) 对全部 \(k\) 成立。

	保持 \(s_k\) 不变时，必须有实际的边界度 \(mk=0\)，即 \(k=0\)。若只保留顶点的值，则可将一胞腔上的映射重新选择为常值，再常值延拓到二胞腔；这对每个 \(k\) 都可做到。这里重新选择不要求与原 \(s_k\) 同伦，所以没有声称非零度的圆周映射能够同伦为常值。

	若取 \(A=X^1\)，相对上链只有二维的 \(\mathbb Z\)，没有一次上链可提供余边界。因而 \(H^2(X,A;\mathbb Z)=\mathbb Z\)，相对首障碍是整数 \(mk\) 本身。相对类为零恰好等价于原 \(s_k\) 能够延拓。相对条件改变的是允许调整的部分，也就改变了应使用的障碍群。
\end{solution}

\begin{exercise}\label{g2:ex:maps-to-sphere-top-dimension}
	设 \(X\) 为维数不超过 \(n\) 的 CW 复形，\(n\ge2\)，\(u\) 为 \(H^n(S^n;\mathbb Z)\) 的正向生成元。证明 \([X,S^n]\to H^n(X;\mathbb Z)\)、\([f]\mapsto f^*u\) 是双射。说明为什么删去维数条件后，\(X=S^3,n=2\) 就给出反例。
\end{exercise}
\begin{solution}
	将映射视为平凡 \(S^n\) 丛的截面。纤维 \(S^n\) 是 \((n-1)\) 连通的简单空间，且已有常值截面；\cref{g2:thm:section-torsor}的维数条件于是满足。以常值截面为原点，差异类给出同伦类与 \(H^n(X;\pi_n(S^n))\cong H^n(X;\mathbb Z)\) 的双射。

	还需识别这张双射。先将 \(f\) 在 \(X^{n-1}\) 上同伦到常值，再用同伦延拓性质得到整个 \(X\) 上的映射。每个正向 \(n\) 胞腔压掉边界后成为 \(S^n\)，差异上链在此处的值就是所得球面映射的度。\(f^*u\) 的胞腔代表在同一胞腔上的值也是这个度，所以两张对应相同。

	若取 \(X=S^3,n=2\)，则 \(H^2(S^3;\mathbb Z)=0\)，但常值映射与 Hopf 映射不是同伦的。事实上 \(S^2\) 单连通，自由同伦不再识别不同的带基点三次类，而 \(\pi_3(S^2)=\mathbb Z\)。更高维胞腔开始看到 \(S^2\) 的高阶同伦群，单个二次上同调类已经不能完成分类。
\end{solution}

\begin{exercise}\label{g2:ex:projective-plane-sphere-map}
	令 \(q:\mathbb RP^2\to\mathbb RP^2/\mathbb RP^1\cong S^2\) 为压掉一骨架的商映射。证明 \(q\) 不是零伦，却在全部正次整数同调上诱导零映射。求 \([\mathbb RP^2,S^2]\) 的元素个数，并用胞腔上链解释这个现象。
\end{exercise}
\begin{solution}
	\(\mathbb RP^2\) 的标准胞腔复形在零、一、二维各有一个胞腔，边界 \(\partial_2\) 为乘 \(2\)、\(\partial_1=0\)。因此 \(H_1(\mathbb RP^2;\mathbb Z)=\mathbb Z/2\)、\(H_2(\mathbb RP^2;\mathbb Z)=0\)。球面的一次同调为零，源空间的二次及更高同调为零，故所有正次 \(q_*\) 都为零。

	对偶上链中 \(\delta^1\) 为乘 \(2\)，于是 \(H^2(\mathbb RP^2;\mathbb Z)=\mathbb Z/2\)。适当选取商球面的定向，\(q\) 把唯一二胞腔以系数 \(1\) 映到球面的二胞腔，所以 \(q^*u\) 由上链 \(1\) 表示，正是模 \(2\) 的非零类。零伦映射的正次上同调拉回为零，因此 \(q\) 不是零伦。

	上一题适用于二维复形与 \(S^2\)，故 \([\mathbb RP^2,S^2]\) 恰有两个元素，分别由常值映射与 \(q\) 代表。此处的二次整数上同调来自一次同调的挠部分，即万有系数定理中的 \(\operatorname{Ext}(\mathbb Z/2,\mathbb Z)\)；它不能只通过二次同调映射读出。
\end{solution}

\begin{exercise}\label{g2:ex:disk-extension-choices}
	在 \(\partial D^n\) 上给定到 \(S^n\) 的常值映射，\(n\ge2\)。求全部延拓 \(D^n\to S^n\) 在固定边界的同伦下的分类，并说明为什么没有存在障碍仍可有无穷多种选择。若改选一个度为 \(r\) 的延拓作为参考，差异坐标怎样改变？
\end{exercise}
\begin{solution}
	常值映射本身提供延拓，所以存在性已经解决。每个延拓将边界送到同一点，故唯一地下降为带基点映射
	\[
		D^n/\partial D^n\cong S^n\longrightarrow S^n.
	\]
	固定边界的同伦也恰好下降为带基点同伦。球面度分类定理因此给出整数分类，任意整数都可由一个相应度的球面映射实现，再复合商映射得到延拓。

	从相对障碍看，\(H^{n+1}(D^n,\partial D^n;\mathbb Z)=0\)，因为相对 CW 对只有一个 \(n\) 胞腔；差异群却是 \(H^n(D^n,\partial D^n;\mathbb Z)=\mathbb Z\)。存在障碍与延拓的差异处在相邻两个次数，前者消失不强迫后者消失。

	以常值延拓为参考时，度为 \(k\) 的延拓坐标为 \(k\)。改用度 \(r\) 的延拓作参考，正负粘合给出的差异为 \(k-r\)。因此全部坐标平移了 \(-r\)，而任意两张延拓的差保持不变。
\end{solution}

\begin{exercise}\label{g2:ex:nonzero-obstruction-group-section}
	设 \(n\ge2\)，考察平凡球丛 \(S^{n+1}\times S^n\to S^{n+1}\)。计算其首障碍群与首障碍类，说明“障碍群不为零”不能推出“没有截面”。再解释 \(H^{n+1}(B;\Pi_n)=0\) 在逐层延拓中实际提供的保证。
\end{exercise}
\begin{solution}
	纤维 \(S^n\) 的首个非零同伦群为 \(\pi_n(S^n)=\mathbb Z\)，平凡丛的运输作用恒等。因此首障碍群为
	\[
		H^{n+1}(S^{n+1};\mathbb Z)=\mathbb Z.
	\]
	但任取 \(v_0\in S^n\)，映射 \(x\mapsto(x,v_0)\) 就是一张全局截面。将它限制到各骨架，所有胞腔边界都已有填充，故障碍上链与障碍类均为零。这是非零障碍群中的零元素。

	障碍群为零时，所有落在这个群中的障碍类都自动为零。于是对于任一当前骨架截面，都能在保持更低一层的条件下调整最高一层，再向下一骨架延拓。这个条件保证整批延拓问题在该步可解；针对一张具体截面，只需它对应的那个类为零。若还有更高维胞腔，则仍需继续处理后续次数。
\end{solution}

\begin{exercise}\label{g2:ex:postnikov-two-stage-maps}
	取正文中由 \(m u^2\)、\(\ell u^2\) 定义的三型 \(E_m,E_\ell\)。设 \(f_a:\mathbb CP^\infty\to\mathbb CP^\infty\) 满足 \(f_a^*u=a u\)，其中 \(a\in\mathbb Z\)。证明：存在覆盖 \(f_a\) 的映射 \(E_m\to E_\ell\)，且它在纤维 \(\pi_3\cong\mathbb Z\) 上诱导乘 \(b\)，当且仅当 \(bm=\ell a^2\)。纤维生成元均按正文的传递约定选取。
\end{exercise}
\begin{solution}
	先设所需映射存在。记两纤维的三次上同调生成元为 \(v_m,v_\ell\)，则纤维映射拉回 \(v_\ell\) 为 \(bv_m\)。传递的自然性给出
	\[
		bm u^2=\mathrm d_4(bv_m)=f_a^*(\mathrm d_4v_\ell)=f_a^*(\ell u^2)=\ell a^2u^2,
	\]
	故条件必要。

	反过来，设 \(bm=\ell a^2\)。取 \(h_b:K(\mathbb Z,4)\to K(\mathbb Z,4)\) 在唯一非零同伦群上诱导乘 \(b\)。映射 \(h_b\kappa_m\) 与 \(\kappa_\ell f_a\) 代表同一个四次上同调类，因此可表性定理给出带基点同伦 \(H:h_b\kappa_m\simeq\kappa_\ell f_a\)。在入向路径模型中，令
	\[
		(b_0,\lambda)\longmapsto\bigl(f_a(b_0),\,\overline{H(b_0,-)}*(h_b\lambda)\bigr).
	\]
	这里横线表示反向走过同伦道路，星号表示依次拼接。新道路从 \(\kappa_\ell f_a(b_0)\) 出发，先到 \(h_b\kappa_m(b_0)\)，再到基点，所以公式确实落在 \(E_\ell\) 中并覆盖 \(f_a\)。在基点纤维上，\(H\) 为常值道路，故诱导的环路映射是 \(\Omega h_b\)，在 \(\pi_3\) 上为乘 \(b\)。这证明条件充分。
\end{solution}

\begin{exercise}\label{g2:ex:postnikov-lift-torsor}
	设 \(n\ge2\)、\(G\) 为阿贝尔群，\(E_\kappa\to B\) 为 \(\kappa:B\to K(G,n+1)\) 的入向路径模型。给定 CW 复形上的映射 \(f:Y\to B\)，并假定它至少有一个提升。证明：在始终覆盖这张固定映射 \(f\) 的同伦关系下，提升类组成 \(H^n(Y;G)\) 的主齐性集，即这个群在提升类上自由且传递地作用。说明为何没有天然指定的零元。
\end{exercise}
\begin{solution}
	提升等价于连续选取道路 \(\lambda_y:\kappa f(y)\to*\)。固定一组选择 \(\lambda_y^0\)。任取连续环路族 \(\omega:Y\to\Omega K(G,n+1)\)，拼接 \(\lambda_y^0*\omega_y\) 给出另一个提升。环路族同伦给出覆盖 \(f\) 的提升同伦，因此这是同伦类上的作用。

	对任意另一提升 \(\lambda_y\)，取基点环路 \(\omega_y=\overline{\lambda_y^0}*\lambda_y\)。插入的往返道路 \(\lambda_y^0*\overline{\lambda_y^0}\) 可保持端点收缩，故 \(\lambda_y^0*\omega_y\) 与 \(\lambda_y\) 在覆盖 \(f\) 的条件下同伦。这证明传递性。如果 \(\lambda^0*\omega\) 与 \(\lambda^0*\omega'\) 在覆盖 \(f\) 的条件下同伦，在每条道路前接上 \(\overline{\lambda^0}\) 并收缩往返部分，就得到 \(\omega\simeq\omega'\)，所以作用自由。

	环路空间的 CW 同伦型为 \(K(G,n)\)，可表性定理把 \([Y,\Omega K(G,n+1)]\) 及环路拼接识别为 \(H^n(Y;G)\) 及其加法。固定 \(\lambda^0\) 后，提升类因而可与这个群一一对应；更换初始提升使该对应整体平移。提升问题本身只给出各提升之间的差，没有指定哪一个应当成为零元。这里固定的是整个基底映射；允许基底映射一起移动会产生另一种分类问题。
\end{solution}

\begin{exercise}\label{g2:ex:postnikov-cp-two-sphere-lift}
	设 \(u\) 为 \(H^2(\mathbb CP^2;\mathbb Z)\) 的标准生成元，\(f_a:\mathbb CP^2\to P_2S^2\simeq\mathbb CP^\infty\) 表示 \(a u\)，其中 \(a\in\mathbb Z\)。求哪些 \(f_a\) 可以提升到 \(P_3S^2\)，哪些可以提升到 \(S^2\)。到 \(S^2\) 的提升允许三角形交换到同伦。用结果说明，只知道目标的二次阶段为何不足以解决四维复形上的映射问题。
\end{exercise}
\begin{solution}
	正文已证明，适当选择纤维生成元后，\(k_4(S^2)=u^2\in H^4(\mathbb CP^\infty;\mathbb Z)\)。因此
	\[
		f_a^*k_4(S^2)=a^2u^2\in H^4(\mathbb CP^2;\mathbb Z).
	\]
	这个四次上同调群由 \(u^2\) 自由生成，所以上式为零当且仅当 \(a=0\)。式\eqref{g2:eq:postnikov-lift-obstruction}遂说明，恰好 \(f_0\) 可以提升到第三阶段。

	若存在提升到 \(S^2\) 的映射，继续复合 \(S^2\to P_3S^2\)，就得到第三阶段的提升，因此仍必须 \(a=0\)。当 \(a=0\) 时，可表性定理说明 \(f_0\) 零伦，故取到 \(S^2\) 的常值映射即可得到交换到同伦的提升。

	二次阶段只检测 \(H^2\)，于是每个整数 \(a\) 都产生一张到 \(P_2S^2\) 的映射。四胞腔还会检测下一阶段的粘合类，它通过平方 \(a^2u^2\) 排除了全部非零 \(a\)。同一结论也可由杯积直接看见：\(S^2\) 的二次类平方为零，任何实际来自 \(S^2\) 的拉回类都必须保留这个关系。
\end{solution}

\begin{exercise}\label{g2:ex:annulus-relative-unit-sections}
	设 \(A=[0,1]\times S^1\)，用 \((r,e^{\mathrm i\theta})\) 表示其点。在平凡定向平面丛 \(A\times\mathbb R^2\) 中，要求单位截面在两个边界上分别为 \(e^{\mathrm i k\theta}\) 和 \(e^{\mathrm i\ell\theta}\)，其中 \(k,\ell\in\mathbb Z\)，两条边界都使用同一个 \(\theta\) 方向。
	\begin{enumerate}
		\item 证明这样的单位截面存在当且仅当 \(k=\ell\)。
		\item 当 \(k=\ell\) 时，证明固定两个边界的截面同伦类恰由整数 \(n\) 分类，并可取代表
		\[
			s_n(r,e^{\mathrm i\theta})=e^{\mathrm i(k\theta+2\pi nr)}.
		\]
	\end{enumerate}
\end{exercise}
\begin{solution}
	1. 单位截面可视为映射 \(f:A\rightarrow S^1\)。两个边界包含映射在 \(A\) 中通过改变 \(r\) 同伦，所以其与 \(f\) 的复合有相同绕数，必要条件为 \(k=\ell\)。反之，\(f(r,e^{\mathrm i\theta})=e^{\mathrm i k\theta}\) 就给出一个延拓。

	2. 给一个满足边界条件的 \(f\)，令 \(g(r,e^{\mathrm i\theta})=f(r,e^{\mathrm i\theta})e^{-\mathrm i k\theta}\)。它在两条边界上恒为 \(1\)。圆环的基本群由内边界生成，故 \(g_*:\pi_1(A)\rightarrow\pi_1(S^1)\) 为零。\cref{g2:prop:covering-lift-criterion}用于 \(t\mapsto e^{2\pi\mathrm it}\)，给出唯一的连续提升 \(G:A\rightarrow\mathbb R\)，使 \(G\) 在内边界恒为 \(0\)。外边界上 \(G\) 为整数值连续函数，因而恒为某个 \(n\in\mathbb Z\)。

	将 \(G\) 线性同伦到 \((r,e^{\mathrm i\theta})\mapsto nr\)，内外边界分别一直保持 \(0,n\)。取指数并乘回 \(e^{\mathrm i k\theta}\)，就得到固定边界的 \(f\simeq s_n\)。若 \(f\) 作固定边界的同伦，则相应提升仍可在内边界规范为零；外边界的整数值在同伦参数下连续，故不能变化。因此不同 \(n\) 对应不同同伦类。这个整数也可沿任何一条由内边界到外边界的径向道路读取：它是消去 \(e^{\mathrm i k\theta}\) 后所得闭路的绕数。
\end{solution}

\begin{exercise}\label{g2:ex:double-mobius-frame}
	令 \(L\rightarrow S^1\) 为 Möbius 线丛。构造 \(L\oplus L\) 的一个整体正交标架，说明一个没有无零截面的线丛与自身直和后可以成为平凡丛，并用纤维的连通性解释这与 \(L\) 的截面障碍为何相容。
\end{exercise}
\begin{solution}
	将 \(L\oplus L\) 写为 \(([0,1]\times\mathbb R^2)/((1,z)\sim(0,-z))\)，Euclid 内积可下降到商。在区间上取
	\[
		v_1(t)=(\cos\pi t,\sin\pi t),\qquad v_2(t)=(-\sin\pi t,\cos\pi t).
	\]
	它们处处正交且长为 \(1\)，并满足 \(v_j(1)=-v_j(0)\)，因此恰好满足商中的端点粘合条件，下降为全局正交标架。映射 \((b,a_1,a_2)\mapsto a_1v_1(b)+a_2v_2(b)\) 给出 \(\varepsilon_{S^1}^2\cong L\oplus L\)。

	单个 \(L\) 的单位纤维是两点集，运输交换这两点，故不能沿一胞腔连续连接端点。直和后单位纤维变成道路连通的圆周，半圈旋转正好连接相反端点，而底空间没有二胞腔来产生下一步障碍。构造的 \(v_1,v_2\) 将这个连通性解释落实成了具体标架。
\end{solution}

\begin{exercise}\label{g2:ex:sphere-one-vector-field-parity}
	对 \(n\ge1\)，证明 \(S^n\) 有处处非零连续切向量场当且仅当 \(n\) 为奇数。偶数维时，从假设存在的单位切向量场构造恒等映射到对径映射的同伦；奇数维时，将环境 Euclid 空间看成复向量空间，直接写出一个光滑单位切向量场。
\end{exercise}
\begin{solution}
	设 \(n\) 为偶数，并假定有无零切向量场。归一化得到 \(s:S^n\rightarrow\mathbb R^{n+1}\)，满足 \(\|s(x)\|=1\) 及 \(\langle s(x),x\rangle=0\)。于是
	\[
		H(x,t)=\cos(\pi t)x+\sin(\pi t)s(x)
	\]
	的范数平方为 \(\cos^2(\pi t)+\sin^2(\pi t)=1\)，所以它是球面内的连续同伦，从 \(H(x,0)=x\) 到 \(H(x,1)=-x\)。\cref{g2:prop:degree-basic,g2:prop:reflection-antipodal-degree}却给出两端的度分别为 \(1\) 与 \((-1)^{n+1}=-1\)，矛盾。

	若 \(n=2q+1\)，将 \(\mathbb R^{2q+2}\) 识别为 \(\mathbb C^{q+1}\)，取 \(s(z)=\mathrm iz\)。复数乘 \(\mathrm i\) 保持实内积范数，并满足 \(\langle\mathrm iz,z\rangle=0\)，故 \(s\) 在单位球面上为光滑单位切向量场。由\cref{g2:prop:sections-orthogonal-reduction}，它还给出 \(\mathrm TS^{2q+1}\cong\varepsilon^1\oplus W\) 的一个分解；这里得到的是一条平凡直和因子。
\end{solution}

\begin{exercise}\label{g2:ex:tautological-plane-circle-sections}
	令 \(L\to\mathbb RP^2\) 为典范实线丛，取带自然内积的秩二丛 \(V=\varepsilon^1\oplus L\)。证明 \(S(V)\) 有截面，并计算其截面在纤维内同伦下的类数。提示：在双覆盖 \(S^2\to\mathbb RP^2\) 上，\(V\) 的纤维坐标变换为 \((a,b)\mapsto(a,-b)\)。
\end{exercise}
\begin{solution}
	平凡因子的单位向量给出两张截面 \(s_+=(1,0)\)、\(s_-=(-1,0)\)，所以存在障碍为零。沿 \(\mathbb RP^2\) 的非平凡基本群元素运输时，单位圆周经反射 \((a,b)\mapsto(a,-b)\)，它在 \(\pi_1(S^1)=\mathbb Z\) 上作用为乘 \(-1\)。因此差异类使用反号整数局部系统 \(\mathbb Z_{\mathrm{sgn}}\)。

	在标准提升胞腔基中，\(\mathbb RP^2\) 的一次边界为 \(g-1\)，二次边界为 \(1+g\)，其中 \(g\) 是覆盖变换。代入反号作用，局部系数上链成为
	\[
		\mathbb Z\xrightarrow{-2}\mathbb Z\xrightarrow{0}\mathbb Z.
	\]
	故 \(H^1(\mathbb RP^2;\mathbb Z_{\mathrm{sgn}})=\mathbb Z/2\)。纤维 \(S^1=K(\mathbb Z,1)\) 没有更高同伦群，\cref{g2:thm:section-torsor}遂给出恰有两个截面同伦类。

	还可辨认两张显式代表。限制到 \(\mathbb RP^1\cong S^1\)，典范线丛成为 Möbius 线丛，单位圆周丛正是\cref{g2:ex:reflection-circle-sections}的反射粘合模型。该题的辐角奇偶性将 \(s_+\) 与 \(s_-\) 分开，因此它们在整个底空间上也不能纤维同伦，恰好代表两类。
\end{solution}

